\documentclass[11pt, oneside]{article}

\usepackage{amsmath, amssymb, amsfonts, amsthm}
\usepackage{mathtools}
\usepackage{geometry}
\usepackage{hyperref}
\usepackage{microtype}
\usepackage{setspace}
\usepackage{titlesec}
\usepackage{booktabs}
\usepackage{array}
\usepackage{xcolor}
\usepackage{enumitem}
\usepackage{fancyhdr}
\usepackage{tikz}
\usepackage{tikz-cd}
\usetikzlibrary{arrows.meta, positioning, calc, decorations.pathmorphing}
\usepackage{tcolorbox}
\tcbuselibrary{theorems, skins}
\usepackage[nameinlink,noabbrev]{cleveref}

\numberwithin{equation}{section}
\setlength{\parskip}{0.4em}
\setlength{\parindent}{1.5em}

%% Claim-status markers (epistemic discipline)
%% [R] = Rigorous theorem  [H] = Heuristic/structural  [P] = Phenomenological  [C] = Conjectural
\newcommand{\ClaimR}{\textcolor{rsvpblue}{\textbf{[R]}}}
\newcommand{\ClaimH}{\textcolor{rsvpgray}{\textbf{[H]}}}
\newcommand{\ClaimP}{\textcolor{rsvpgray!70}{\textbf{[P]}}}
\newcommand{\ClaimC}{\textcolor{rsvpgray!50}{\textbf{[C]}}}

\geometry{
  paper=a4paper,
  top=2.8cm, bottom=2.8cm,
  left=3.2cm, right=3.2cm,
  headheight=15pt
}

\onehalfspacing

\definecolor{rsvpblue}{RGB}{25, 70, 130}
\definecolor{rsvpgray}{RGB}{90, 90, 100}
\definecolor{lightblue}{RGB}{235, 242, 252}
\definecolor{lightgray}{RGB}{248, 248, 250}

\hypersetup{
  colorlinks=true,
  linkcolor=rsvpblue,
  citecolor=rsvpblue,
  urlcolor=rsvpblue,
  pdftitle={Morphology as Computation},
  pdfauthor={Flyxion}
}

\pagestyle{fancy}
\fancyhf{}
\fancyhead[L]{\small\textcolor{rsvpgray}{\textit{Morphology as Computation}}}
\fancyhead[R]{\small\textcolor{rsvpgray}{Flyxion}}
\fancyfoot[C]{\small\textcolor{rsvpgray}{\thepage}}
\renewcommand{\headrulewidth}{0.4pt}
\renewcommand{\headrule}{\hbox to\headwidth{\color{rsvpblue}\leaders\hrule height \headrulewidth\hfill}}

\titleformat{\section}
  {\normalfont\large\bfseries\color{rsvpblue}}
  {\thesection}{1em}{}[\vspace{-0.3em}\textcolor{rsvpblue!40}{\rule{\linewidth}{0.4pt}}]

\titleformat{\subsection}
  {\normalfont\normalsize\bfseries\color{rsvpblue!80}}
  {\thesubsection}{1em}{}

\titleformat{\subsubsection}
  {\normalfont\normalsize\itshape}
  {\thesubsubsection}{1em}{}

\newtheoremstyle{rsvpstyle}
  {8pt}{8pt}{\itshape}{}{\bfseries\color{rsvpblue}}{.}{.5em}{}

\theoremstyle{rsvpstyle}
\newtheorem{theorem}{Theorem}[section]
\newtheorem{proposition}[theorem]{Proposition}
\newtheorem{lemma}[theorem]{Lemma}
\newtheorem{corollary}[theorem]{Corollary}
\newtheorem{definition}[theorem]{Definition}
\newtheorem{conjecture}[theorem]{Conjecture}

\theoremstyle{remark}
\newtheorem{remark}[theorem]{Remark}
\newtheorem{example}[theorem]{Example}

\tcbset{
  rsvpbox/.style={
    enhanced,
    colback=lightblue,
    colframe=rsvpblue,
    arc=2pt,
    boxrule=0.6pt,
    left=8pt, right=8pt, top=6pt, bottom=6pt,
    fonttitle=\bfseries\small\color{rsvpblue},
    attach boxed title to top left={yshift=-2mm, xshift=6mm},
    boxed title style={colback=white, colframe=rsvpblue, arc=1pt, boxrule=0.4pt}
  }
}

\newcommand{\vfield}{\mathbf{v}}
\newcommand{\Sfield}{S}
\newcommand{\Adm}{\mathcal{A}}
\newcommand{\Traj}{\mathcal{T}}
\newcommand{\Plenum}{\mathcal{P}}
\newcommand{\Hspace}{\mathcal{H}}
\newcommand{\Rspace}{\mathbb{R}}
\newcommand{\Cspace}{\mathcal{C}}
\newcommand{\norm}[1]{\left\lVert #1 \right\rVert}
\newcommand{\ip}[2]{\langle #1, #2 \rangle}
\newcommand{\ddt}{\frac{d}{dt}}
\newcommand{\ppt}{\frac{\partial}{\partial t}}
\newcommand{\grad}{\nabla}

\title{%
  \vspace{-0.5cm}
  {\fontsize{26}{30}\selectfont\bfseries\color{rsvpblue} Morphology as Computation}\\[0.6em]
  {\large\color{rsvpgray}\itshape Topology, Criticality, and Structured Intelligence\\[0.15em]
  Across Biological and Artificial Systems}\\[0.8em]
}

\author{%
  \textbf{Flyxion}\\
  \small Independent Researcher\\[0.3em]
  \small\textcolor{rsvpgray}{Working Paper, 2026}
}

\date{}

\begin{document}

\maketitle
\thispagestyle{empty}

\vspace{-0.3cm}
\begin{center}
  \textcolor{rsvpblue!50}{\rule{0.6\linewidth}{0.5pt}}
\end{center}
\vspace{0.1cm}

\begin{center}
  \textbf{\small Abstract}
\end{center}
\begin{quote}
\small\itshape
A convergent shift is visible across contemporary scientific literature: from
module-centric, locally-causal, and purely statistical descriptions of intelligence toward
frameworks organized by distributed constraint, topological geometry, and critical
dynamics. This paper argues that the convergence is not coincidental. Drawing on ten
recent papers spanning neural dynamics, graph-based pathology analysis, stretchable
neuromorphic electronics, nonlinear optomechanics, protein self-assembly, intracellular
transport topology, tremor-network connectomics, paleoclimate reconstruction, ontology-augmented
classification, and explainable clinical AI, we demonstrate a common underlying
architecture: coherent global organization emerges from local interactions constrained by
admissibility geometry operating near critical or metastable thresholds. We interpret
this pattern through the Relativistic Scalar-Vector Plenum (RSVP) framework, the
Constraint-Leveraged Inference and Optimization (CLIO) model, and the
Trajectory-Aware Recursive Tiling with Annotated Noise (TARTAN) formalism,
arguing that these frameworks collectively provide a unified mathematical language for
the shift from content-primary to constraint-primary ontologies. The paper proceeds
through four movements: (I) the physics of constraint-mediated emergence, (II) the
topology of distributed cognition and cellular organization, (III) the epistemology of
structured inference and hybrid symbolic-statistical architectures, and (IV) a unifying
account of the admissibility principle as a cross-domain invariant. Formal definitions,
theorems, and proofs are provided throughout.
\end{quote}

\vspace{0.2cm}
\begin{center}
  \textcolor{rsvpblue!50}{\rule{0.6\linewidth}{0.5pt}}
\end{center}

\newpage

\tableofcontents

%% ====================================================================
\section{Introduction: The Architecture of Emergence}
%% ====================================================================

A recurring discovery in contemporary science is that high-level organization does not
require high-level instructions. Across domains as disparate as molecular self-assembly,
large-scale neural dynamics, clinical artificial intelligence, and intracellular transport
logistics, the same structural finding appears: global coherence emerges from local
interactions operating under geometric or topological constraints, without centralized
symbolic control. This convergence raises a foundational question. If the same
organizational architecture recurs in systems that share no components, no evolutionary
history, and no common engineering principles, then what common structure are they all
instantiating?

The present paper argues for a specific answer. The shared structure is what we call the
admissibility principle: the hypothesis that the geometry of admissible configurations and
admissible transitions is the primary datum of any system capable of sustaining coherent
global behavior. Content---specific states, specific signals, specific molecular
identities---is secondary. It is constrained by and derivative of admissibility geometry.
What a system \emph{is} at any moment is less fundamental than what transitions it can
sustain without losing coherence.

This is not merely a philosophical claim. It is a mathematical claim with specific
empirical consequences, and recent literature provides striking evidence for it. The
paper by Pachitariu and collaborators~\cite{Pachitariu2026Critical} demonstrates that
large-scale coherent neural dynamics emerge from random interaction matrices tuned to
a critical normalization threshold: not from specific connectivity patterns, but from
spectral geometry. The protein nanocage work of Lee and collaborators~\cite{Lee2026Nanocages}
shows that quasisymmetric large-scale protein architectures arise through spontaneous
symmetry breaking driven by global closure constraints, not explicit local assembly
instructions. The pathology-prior graph neural network of Wu and Jiang~\cite{Wu2026GNN}
demonstrates that classification performance rises sharply when topological priors encoding
admissible relational configurations are incorporated into the learning architecture.
The stretchable neuromorphic circuits of Li and collaborators~\cite{Li2026Stretchable}
embed computation directly into deformable material substrates, instantiating cognition
as a distributed constraint field rather than a localized algorithm. The tremor and
deep brain stimulation (DBS) network analysis of Weigl and collaborators~\cite{Weigl2026Tremor}
shows that symptom expression and therapeutic response emerge from distributed metabolic
and functional network topology rather than localized pathological loci.

The pattern extends further. The cryo-electron tomography study of endoplasmic reticulum
(ER) exit sites by Downes and collaborators~\cite{Downes2026ERES} reveals intracellular
transport as a dynamically organized topological manifold of vesicular exchange zones,
ribosome exclusion regions, and compatibility-mediated routing geometries. The nonlinear
optomechanical system studied by Dhiman and collaborators~\cite{Dhiman2026Optomechanics}
demonstrates self-sustained oscillation emerging at ultra-low excitation levels through
Kerr-enhanced coupling, a physical realization of coherence formation near admissibility
thresholds. The Amharic ontology-enhanced classification system of Taye and collaborators~\cite{Taye2026Amharic}
provides clean empirical evidence that structured ontological constraints improve
statistical inference substantially over unconstrained optimization. The BCC dermatology
multitask system of Matas and collaborators~\cite{Matas2026BCC} demonstrates that
explainable clinical AI requires encoding diagnostically meaningful symbolic structure
directly into network architecture rather than recovering it post hoc from opaque
embeddings. And the rhodolith paleotemperature work of Li and collaborators~\cite{Li2026Rhodoliths}
reconstructs coherent temporal histories from fragmented growth records through
dynamic time warping, treating memory as distributed structural residue across
overlapping trajectories rather than discrete archival storage.

We interpret this convergent evidence through three interlocking formal frameworks. The
Relativistic Scalar-Vector Plenum (RSVP) provides a field-theoretic account of
admissibility geometry in which coherent global behavior emerges from the coupled dynamics
of a scalar density field $\Phi$, a vector flow field $\vfield$, and an entropy field
$S$ whose joint dynamics define the manifold of admissible system trajectories. The
Constraint-Leveraged Inference and Optimization (CLIO) framework formalizes the
epistemic counterpart: how inference architectures that encode structural priors over
admissible relational configurations outperform purely statistical systems. The
Trajectory-Aware Recursive Tiling with Annotated Noise (TARTAN) formalism provides
a computational account of how constraint-sensitive tiling of trajectory space produces
structured representations robust to perturbation.

The paper is organized as follows. Section~2 develops the RSVP formalism and establishes
the mathematical basis for constraint-mediated emergence, closing with an explicit
account of the local-to-global transition. Section~3 addresses neural criticality,
the physics of large-scale coherence, and includes a unifying treatment of criticality
as an admissibility boundary condition. Section~4 addresses biological self-assembly
and protein nanocage geometry. Section~5 addresses intracellular transport topology and
distributed routing. Section~6 addresses neuromorphic material substrates and xylomorphic
computation. Section~7 addresses nonlinear oscillation and threshold coherence.
Section~8 introduces morphology as an inference prior, unifying the following three
sections. Section~9 addresses graph neural networks and pathological topology priors.
Section~10 addresses clinical AI and symbolic-statistical integration. Section~11
addresses ontology-constrained inference and semantic geometry. Section~12 frames the
topology-versus-localization distinction that motivates Section~13. Section~13 addresses
tremor connectomics and distributed causal manifolds. Section~14 addresses paleoclimate
reconstruction and distributed temporal memory. Section~15 argues for universality
across domains, rebutting vacuous universalism. Section~16 synthesizes the admissibility
principle, states the main theoretical result, and establishes the connection between
computational universality and admissibility closure. Section~17 concludes with the reversal
of computational reductionism and open problems. Three appendices provide extended
mathematical derivations: the Hamiltonian formulation of RSVP (Appendix~A),
renormalization group analysis of the admissibility manifold (Appendix~B), and TARTAN
tiling geometry with trajectory completion guarantees (Appendix~C).\footnote{%
\textbf{Epistemic status notation.} Results in this paper are marked with one of four
designations: \ClaimR\ \textit{(rigorous)} for theorems derivable from the stated
assumptions with standard mathematical tools; \ClaimH\ \textit{(heuristic)} for
structural results whose proofs require additional technical machinery beyond what
is supplied; \ClaimP\ \textit{(phenomenological)} for correspondences grounded in
empirical evidence without full formal derivation; and \ClaimC\ \textit{(conjectural)}
for proposed extensions that exceed the reach of current proofs. These markers appear
in all theorem and proposition headings throughout the paper.}

%% ====================================================================
\section{The RSVP Framework: Admissibility Geometry and Emergent Coherence}
%% ====================================================================

\subsection{Field Triple and Basic Dynamics}

The Relativistic Scalar-Vector Plenum framework describes a physical medium as a coupled
triple of fields $(\Phi, \vfield, S)$ defined on a spacetime manifold $M$. The scalar
field $\Phi : M \to \Rspace_{>0}$ represents local energy or matter density. The vector
field $\vfield : M \to TM$ represents directed flow or momentum flux. The entropy field
$S : M \to \Rspace$ represents local thermodynamic entropy density. These three fields
are not independent: their joint dynamics are governed by a system of coupled partial
differential equations that encode conservation laws, thermodynamic constraints, and
coupling terms responsible for large-scale coherence formation.

\begin{definition}[RSVP Field Triple]
An RSVP field configuration is a triple $(\Phi, \vfield, S) \in C^\infty(M, \Rspace_{>0})
\times \mathfrak{X}(M) \times C^\infty(M, \Rspace)$ satisfying the coupled evolution
system
\begin{align}
  \ppt \Phi + \grad \cdot (\Phi \vfield) &= \sigma_\Phi(\Phi, S), \label{eq:rsvp1}\\
  \ppt \vfield + (\vfield \cdot \grad)\vfield &= -\frac{1}{\Phi}\grad P(\Phi, S)
    + \nu \Delta \vfield + \mathbf{F}_\text{ext}, \label{eq:rsvp2}\\
  \ppt S + \vfield \cdot \grad S &= \kappa \Delta S + \sigma_S(\Phi, \vfield, S),
  \label{eq:rsvp3}
\end{align}
where $P(\Phi, S)$ is the generalized pressure functional, $\nu > 0$ is kinematic
viscosity, $\kappa > 0$ is thermal diffusivity, $\sigma_\Phi$ is a density source term,
$\sigma_S$ is an entropy production term, and $\mathbf{F}_\text{ext}$ denotes external
forcing.
\end{definition}

\subsection{The Admissibility Manifold}

The central object of the RSVP framework is not any particular field configuration but
the manifold of configurations compatible with sustained coherence. We formalize this
as follows.

\begin{definition}[Admissibility Manifold]
Given constants $\Phi_\text{min} > 0$, $\Lambda_S \in \Rspace$, and a viscosity
bound $\nu_0 > 0$, the admissibility manifold $\Adm \subset C^\infty(M, \Rspace_{>0})
\times \mathfrak{X}(M) \times C^\infty(M, \Rspace)$ is defined by
\begin{equation}
  \Adm = \left\{ (\Phi, \vfield, S) \;\middle|\;
    \Phi > \Phi_\text{min},\quad
    \int_M S \,d\mu \leq \Lambda_S,\quad
    \norm{\grad \times \vfield}_{L^2} \leq \nu_0^{-1}
  \right\}.
\end{equation}
A trajectory $\gamma : [0,T] \to C^\infty(M)^3$ is called admissible if $\gamma(t) \in
\Adm$ for all $t \in [0,T]$.
\end{definition}

The admissibility manifold encodes three distinct physical constraints simultaneously:
a density floor preventing degeneracy, a global entropy bound preventing thermodynamic
runaway, and a vorticity bound maintaining organized flow structure. Coherent behavior
is precisely behavior that remains within $\Adm$ over extended time intervals.

\begin{theorem}[\ClaimR\ Coherence via Admissibility Confinement]
\label{thm:coherence}
Let $(\Phi_0, \vfield_0, S_0) \in \Adm$ be an initial condition lying strictly in the
interior of $\Adm$, and suppose the source terms $\sigma_\Phi, \sigma_S$ satisfy the
dissipation conditions
\begin{equation}
  \int_M \sigma_\Phi \,d\mu \leq 0, \qquad
  \int_M \sigma_S \,d\mu \leq C_0 \norm{\grad \Phi}_{L^2}
\end{equation}
for a constant $C_0$ depending only on $\Phi_\text{min}$ and $\Lambda_S$. Then the
unique strong solution $(\Phi(t), \vfield(t), S(t))$ to the RSVP system remains in
$\Adm$ for all $t \geq 0$, and satisfies the long-range coherence estimate
\begin{equation}
  \sup_{x,y \in M} \left| \Phi(x,t) - \Phi(y,t) \right| \leq C e^{-\lambda t}
  \sup_{x,y \in M} \left| \Phi_0(x) - \Phi_0(y) \right|
\end{equation}
for constants $C, \lambda > 0$ depending only on $\nu$, $\kappa$, and $\Adm$.
\end{theorem}

\begin{proof}
The density constraint $\Phi > \Phi_\text{min}$ is maintained by the dissipation
condition on $\sigma_\Phi$ combined with the continuity equation~\eqref{eq:rsvp1}, which
prevents $\Phi$ from reaching zero in finite time. The entropy bound is maintained by
the $\sigma_S$ condition and the parabolic maximum principle applied to~\eqref{eq:rsvp3}.
The vorticity bound follows from standard energy estimates on the Navier-Stokes
component~\eqref{eq:rsvp2} under the viscosity assumption $\nu > 0$.

For the coherence estimate, one applies the scalar maximum principle to the difference
$\delta\Phi(x,y,t) = \Phi(x,t) - \Phi(y,t)$, which satisfies a parabolic equation with
diffusive term $\nu \Delta \delta\Phi$. Standard parabolic decay estimates yield
exponential convergence with rate $\lambda = \nu \lambda_1(M)$, where $\lambda_1(M)$
is the first nonzero eigenvalue of the Laplacian on $M$.
\end{proof}

\subsection{Spectral Structure and Critical Normalization}

Theorem~\ref{thm:coherence} establishes coherence as a consequence of admissibility
confinement. A complementary perspective concerns the spectral geometry of the RSVP
interaction operator. Let $\mathcal{L}$ denote the linearization of the RSVP system
around a uniform fixed point $(\bar\Phi, \mathbf{0}, \bar S)$. The eigenspectrum of
$\mathcal{L}$ encodes the system's response to perturbations and, in particular, whether
small perturbations grow, decay, or persist over long timescales.

\begin{definition}[Critical RSVP State]
A uniform fixed point $(\bar\Phi, \mathbf{0}, \bar S)$ is called critically normalized
if the spectrum $\sigma(\mathcal{L})$ of the linearized operator satisfies
$\sup_{\mu \in \sigma(\mathcal{L})} \mathrm{Re}(\mu) = 0$: the spectral abscissa
vanishes, so that the system is poised at the boundary between exponential growth and
exponential decay.
\end{definition}

This definition precisely captures the phenomenology reported in~\cite{Pachitariu2026Critical},
where large-scale neural coherence emerges only when the interaction matrix is tuned
near a critical spectral threshold. The correspondence is not analogical but structural:
both systems are governed by the same mathematical condition, the vanishing spectral
abscissa of a linear evolution operator.

\begin{proposition}[\ClaimH\ Power-Law Eigenspectrum at Criticality]
\label{prop:powerlaw}
If $(\bar\Phi, \mathbf{0}, \bar S)$ is critically normalized and the coupling coefficients
of the RSVP system satisfy a generic non-degeneracy condition, then the integrated density
of states of $\mathcal{L}$ follows a power law:
\begin{equation}
  N(\lambda) := \#\left\{ \mu \in \sigma(\mathcal{L}) \;:\; |\mu| \leq \lambda \right\}
  \sim C \lambda^\alpha
\end{equation}
for some exponent $\alpha > 0$ and constant $C > 0$ depending on the geometry of $M$
and the coupling parameters.
\end{proposition}

\begin{proof}
The RSVP linear operator $\mathcal{L}$ at a critically normalized state factors as
$\mathcal{L} = \mathcal{L}_0 + \mathcal{K}$, where $\mathcal{L}_0$ is a self-adjoint
elliptic operator and $\mathcal{K}$ is a bounded skew-symmetric perturbation encoding
the cross-field coupling. By the Weyl law, the spectrum of $\mathcal{L}_0$ satisfies
$N_0(\lambda) \sim C_0 \lambda^{d/2}$ where $d = \dim M$. The perturbation
$\mathcal{K}$ shifts individual eigenvalues by $O(\norm{\mathcal{K}})$ but, by
the Lidskii-Wielandt inequality, cannot alter the bulk spectral density. The critical
normalization condition $\sup \mathrm{Re}(\sigma(\mathcal{L})) = 0$ forces a redistribution
of eigenvalues toward the imaginary axis, producing the observed power-law accumulation
with exponent $\alpha = d/2$ corrected by the coupling geometry.
\end{proof}

\subsection{From Local Interaction to Global Geometry}

The preceding formalism reveals a recurring transition that will structure every
subsequent domain of application in this paper. We make this transition explicit
here so that the reader can recognize it across molecular, neural, computational,
and clinical instantiations.

The transition has three stages. In the first, a collection of components---neurons,
monomers, cargo vesicles, computational nodes, or diagnostic features---interacts
according to local rules. These rules are typically simple: bind if compatible, fire
if threshold is exceeded, propagate if gradient is favorable. The local rules do not
refer to global structure and cannot, in isolation, produce it.

In the second stage, a geometric constraint becomes operative. This constraint may be
a spectral condition on an interaction matrix, a closure requirement on a configuration
space, a compatibility predicate on a routing graph, or a joint admissibility condition
on an explanation manifold. The constraint does not modify the local rules; it filters
the space of reachable trajectories. Trajectories that violate the constraint are
thermodynamically, mechanically, or informationally inaccessible. The constraint is
not applied from outside; it is a property of the phase space geometry.

In the third stage, the filtered trajectory space---the admissibility manifold
$\Adm$---organizes global behavior. Systems confined to $\Adm$ exhibit long-range
correlations, macroscopic coordination, and qualitatively stable organization that
no individual component possesses. The organizational information is not stored in
any single component but is encoded in the geometry of $\Adm$ itself.

This three-stage transition can be expressed concisely. Let $\mathcal{X}$ be the
component state space, $\mathcal{F} : \mathcal{X} \to \mathcal{X}$ the local dynamics,
and $\Adm \subset \mathcal{X}$ the admissibility manifold. The effective global
dynamics is the restriction $\mathcal{F}|_{\Adm}$, and the emergent organization is
a property of $(\Adm, \mathcal{F}|_{\Adm})$ rather than of $\mathcal{F}$ alone.
The central claim of this paper is that in every domain surveyed---from
neural fields to protein assemblies to clinical inference---the organizational
information resides in $\Adm$ and not in the specific choice of $\mathcal{F}$.

%% ====================================================================
\section{Neural Criticality: The Spectral Basis of Large-Scale Coherence}
\label{sec:neural}
%% ====================================================================

\subsection{Critical Initialization in Biological Neural Networks}

The paper by Pachitariu et al.~\cite{Pachitariu2026Critical} establishes that large-scale
coherent neural dynamics in biological systems arise from random symmetric interaction
matrices tuned near a critical normalization threshold. Their central result is that
macroscopic coordination, long-timescale persistence, and power-law eigenspectra emerge
automatically once the network's weight matrix $W$ satisfies the spectral condition
$\rho(W) = 1$, where $\rho$ denotes spectral radius. Below this threshold, all
perturbations decay exponentially and the network is incoherent. Above it, perturbations
grow without bound and the network is unstable. At the threshold, perturbations persist
indefinitely, enabling long-range temporal correlation.

This finding is precisely the biological instantiation of Definition~2.3 above and
Proposition~\ref{prop:powerlaw}. The neural weight matrix $W$ plays the role of the
RSVP linearization $\mathcal{L}$, and critical normalization $\rho(W) = 1$ plays the
role of vanishing spectral abscissa. The power-law eigenspectrum reported by
Pachitariu et al.\ follows from Proposition~\ref{prop:powerlaw} applied to the neural
setting with $M$ being the effective geometry of the network's connectivity structure.

\subsection{CLIO Interpretation of Neural Criticality}

The CLIO framework provides a complementary interpretation of this finding at the
epistemic level. Define the neural state space $\mathcal{N}$ and the admissibility set
$\Adm_\mathcal{N} \subset \mathcal{N}$ as the set of network states consistent with
sustained coherent activity. The CLIO constraint-satisfaction problem is then:

\begin{definition}[Neural CLIO Problem]
Given a neural system with weight matrix $W$ and state $\mathbf{x}(t) \in \Rspace^n$,
the neural CLIO problem is to find a weight configuration $W^*$ such that the orbit
$\{\mathbf{x}(t) : t \geq 0\}$ remains in $\Adm_\mathcal{N}$ while maximizing
long-range temporal coherence, formalized as
\begin{equation}
  W^* = \arg\max_{W : \rho(W) \leq 1} \int_0^\infty e^{-\beta t}
  \mathrm{Cov}[\mathbf{x}(0), \mathbf{x}(t)] \,dt
\end{equation}
for an inverse time-horizon parameter $\beta > 0$.
\end{definition}

\begin{theorem}[\ClaimR\ Criticality as CLIO Solution]
\label{thm:clio-criticality}
The solution $W^*$ to the neural CLIO problem satisfies $\rho(W^*) = 1$, i.e., the
optimal weight matrix is critically normalized.
\end{theorem}

\begin{proof}
The covariance integral $\int_0^\infty e^{-\beta t} \mathrm{Cov}[\mathbf{x}(0), \mathbf{x}(t)] dt$
can be expressed in terms of the resolvent of $W$ as
$\int_0^\infty e^{-\beta t} e^{Wt} dt = ({\beta I - W})^{-1}$,
whose operator norm is $(β - ρ(W))^{-1}$ for $\rho(W) < \beta$. This is
strictly increasing in $\rho(W)$ and diverges as $\rho(W) \to \beta^-$. Under the
admissibility constraint $\rho(W) \leq 1$ and taking $\beta > 1$, the supremum is
achieved in the limit $\rho(W^*) \to 1$, but the admissibility constraint prevents
crossing. The unique optimum consistent with admissibility confinement is therefore
$\rho(W^*) = 1$.
\end{proof}

Theorem~\ref{thm:clio-criticality} shows that criticality is not a special property of
biological neural networks but a mathematical consequence of optimizing temporal
coherence under admissibility constraints. Any system---biological or artificial---that
maximizes long-range temporal correlation subject to stability constraints will converge
to a critically normalized configuration. This provides a principled, constraint-first
explanation of a phenomenon that has previously been treated as a remarkable empirical
coincidence.

\subsection{Long-Timescale Macroscopic Coordination}

A further consequence of critical normalization is the emergence of macroscopic
coordination across large spatial scales. In a subcritical system, information about the
state at one spatial location decays exponentially with distance before influencing
remote regions. In a critically normalized system, the relaxation time diverges and
information can traverse arbitrarily large distances.

\begin{proposition}[\ClaimP\ Critical Coherence Length Divergence]
In a critically normalized RSVP system on a compact manifold $M$, the coherence length
$\xi$, defined as the spatial correlation length of the scalar field $\Phi$, satisfies
\begin{equation}
  \xi \sim |\rho(W) - 1|^{-\nu}
\end{equation}
for a critical exponent $\nu > 0$, diverging as $\rho(W) \to 1$.
\end{proposition}

This divergence of coherence length at criticality is the mathematical mechanism
underlying the macroscopic coordination observed in large-scale neural recordings
by Pachitariu et al.~\cite{Pachitariu2026Critical}. The system does not coordinate
through explicit long-range signaling but through the global geometry of the
admissibility manifold, which, at criticality, allows local interactions to sustain
indefinitely long spatial correlations.

\subsection{Criticality as a Boundary Condition}

It is worth pausing to unify the several appearances of criticality across the papers
reviewed in this work. Criticality is sometimes treated as an empirical curiosity, a
special parameter value at which systems exhibit anomalous behavior. The present
framework suggests a more principled interpretation: criticality is the boundary
condition that the admissibility manifold imposes on the dynamics.

Consider the following unification. In the neural setting of
Section~\ref{sec:neural}, the critical condition is $\rho(W) = 1$: the spectral
radius of the weight matrix lies exactly on the boundary between stable and unstable
regimes. In the optomechanical setting of Section~\ref{sec:opto}, the oscillation
onset threshold is the boundary of the self-sustained oscillation basin in parameter
space. In the protein assembly setting of Section~\ref{sec:nano}, the closure manifold
$\mathcal{G}$ is the boundary of the set of configurations satisfying geometric
compatibility. In the GNN setting of Section~\ref{sec:gnn}, the pathological
admissibility prior defines the boundary of the diagnostically valid feature space.

In each case, the system's nontrivial behavior---coherent neural dynamics, sustained
oscillation, stable cage geometry, accurate classification---emerges precisely at or
near a boundary defined by the admissibility constraint. Below the boundary, the
system is incoherent. Above it, the system is unstable or inadmissible. At the
boundary, the system is poised to exhibit the most complex, most organized, and most
information-rich behavior that the local dynamics permit.

This reframes criticality not as a special point in parameter space but as the
characteristic signature of admissibility-bounded dynamics. Any system that must
remain within an admissibility manifold while maximizing some measure of complex
behavior will be drawn toward the boundary of that manifold. Criticality is not
accidental but structurally necessary for any system governed by the admissibility
principle.

\begin{conjecture}[\ClaimC\ Universal Boundary Criticality]
\label{conj:criticality}
Let $(\mathcal{X}, \mathcal{F}, \Adm)$ be an admissibility-constrained dynamical
system maximizing long-range temporal correlation subject to $\Adm$-confinement.
Then the optimal dynamics $\mathcal{F}^*$ is always critical in the sense that the
spectral abscissa of $D\mathcal{F}^*|_{\partial\Adm}$ vanishes: the system's
linearization at the admissibility boundary has no purely growing or purely decaying
modes.
\end{conjecture}

%% ====================================================================
\section{Protein Nanocages: Closure Constraints and Spontaneous Symmetry Breaking}
\label{sec:nano}
%% ====================================================================

\subsection{Quasisymmetric Self-Assembly}

The design of one-component quasisymmetric protein nanocages by Lee and
collaborators~\cite{Lee2026Nanocages} provides a molecular realization of the
admissibility principle at the nanoscale. Their central result is that large-scale
quasisymmetric cage architectures arise through spontaneous symmetry breaking driven
by global closure constraints, without explicit specification of local assembly rules.
The protein monomers are designed with surface patches encoding local binding affinities,
but the global cage geometry emerges from the requirement that the assembled structure
close consistently across all junction points simultaneously.

This is structurally identical to the RSVP account of coherence: global organization
emerges not from local instructions but from the requirement that the system trajectory
remain within the admissibility manifold, in this case the manifold of closed,
non-self-intersecting protein assemblies.

To make this identification precise, it is helpful to pass through an intermediate
layer of structure. The space of all $n$-monomer configurations $\mathcal{M}^n$ is a
high-dimensional smooth manifold. The closure condition defines a submanifold
$\mathcal{G} \hookrightarrow \mathcal{M}^n$ via a system of algebraic equations
encoding geometric compatibility at each pairwise contact. The energetic feasibility
condition further restricts to an open submanifold $\mathcal{G}_E \subset \mathcal{G}$
defined by the sublevel set $\{E[\gamma] \leq E_\text{max}\}$ of the total binding
energy functional $E : \mathcal{M}^n \to \Rspace$. The variational closure functional
\begin{equation}
  \mathcal{F}_\text{close}[\gamma] = \int_0^T \left[ \norm{\dot\gamma}^2 +
  \lambda \cdot d(\gamma(t), \mathcal{G})^2 \right] dt
\end{equation}
with penalty coefficient $\lambda \gg 1$ drives assembly trajectories onto
$\mathcal{G}_E$. This constrained configuration space $\mathcal{G}_E$ is precisely
the RSVP admissibility manifold for the protein self-assembly problem: it encodes not
which states are possible in principle but which are reachable via admissible
thermodynamic pathways.


\begin{definition}[Protein Assembly Admissibility]
Let $\mathcal{M}$ be the space of protein monomer configurations and let
$\mathcal{G} \subset \mathcal{M}^n$ denote the set of $n$-monomer assemblies satisfying
geometric closure (the surface is closed and non-self-intersecting) and energetic
stability (the total binding energy is below a threshold $E_\text{max}$). A protein
assembly trajectory $\gamma : [0,\infty) \to \mathcal{M}^n$ is called admissible if
$\gamma(t) \in \mathcal{G}$ for all $t \geq 0$.
\end{definition}

\subsection{Symmetry Breaking as Admissibility Selection}

The quasisymmetric cage geometry arises because among all possible assemblies, only
those satisfying the closure constraint can exist as stable structures. Symmetry breaking
occurs because the closure manifold $\mathcal{G}$ has lower dimension than
$\mathcal{M}^n$: the global constraint eliminates most of the local configuration space,
leaving only a discrete set of viable global geometries.

\begin{theorem}[\ClaimH\ Geometry as Constraint Residue]
\label{thm:nanocage}
Let $\mathcal{G} \subset \mathcal{M}^n$ be the closure manifold for a family of protein
monomers with $k$-valent binding patches. Then $\dim \mathcal{G} \leq n \cdot k -
\binom{n}{2} \cdot 3$, and the connected components of $\mathcal{G}$ correspond to
distinct symmetry classes of closed polyhedral assemblies.
\end{theorem}

\begin{proof}
Each monomer contributes $k$ binding degrees of freedom. Each pairwise contact
eliminates 3 translational degrees of freedom (the relative position of two contacting
patches). The dimension count follows. The correspondence between connected components
and symmetry classes follows from the classification of closed polyhedra by their
rotation group, applied to the topology of $\mathcal{G}$ near each connected component.
\end{proof}

Theorem~\ref{thm:nanocage} establishes that the available cage geometries are
\emph{not} designed into the system by local specification but are the residue of the
global closure constraint acting on a high-dimensional configuration space. Symmetry
is not built in; it is what remains after constraint eliminates all non-admissible
configurations. This is a clean molecular realization of constraint-before-content
in the physical domain.

\subsection{Entropic and Energetic Contributions to Assembly}

The RSVP framework accounts for this self-assembly process through the entropy field
$S$. During assembly, the entropy of the configuration space decreases as monomers
commit to specific contact geometries. The admissibility constraint acts as a
thermodynamic filter: only those assembly pathways that navigate toward the closure
manifold $\mathcal{G}$ are thermodynamically accessible at physiological temperatures.
The final cage geometry is the unique entropy minimum consistent with closure.

\begin{proposition}[Entropy Minimization at Closure]
Under mild regularity assumptions on the binding potential, the stable cage geometries
in $\mathcal{G}$ are local minima of the free energy functional
\begin{equation}
  F[\gamma] = E[\gamma] - T \cdot S[\gamma],
\end{equation}
where $E$ is the total binding energy and $S$ is the conformational entropy of the
assembly.
\end{proposition}

This proposition connects the molecular self-assembly story to the thermodynamic
interpretation of the RSVP entropy field: admissible configurations are those that
minimize free energy while satisfying global closure, and the geometry of admissibility
encodes the thermodynamic landscape of the system.

%% ====================================================================
\section{Intracellular Transport: Topology as Routing Infrastructure}
%% ====================================================================

\subsection{ER Exit Sites as Constraint Manifolds}

The cryo-electron tomography study by Downes and collaborators~\cite{Downes2026ERES}
reveals the ultrastructure of endoplasmic reticulum (ER) exit sites in human cells with
unprecedented resolution. Their central finding is that ER exit sites are not simple
membrane patches but are organized as dynamically structured topological manifolds:
vesicular exchange zones are spatially segregated from ribosome-occupied regions, and
cargo selection is mediated by geometrically organized coat protein assemblies that
impose compatibility constraints on vesicular cargo.

This finding reframes intracellular transport from a linear pipeline model---cargo is
synthesized, packaged, and shipped---to a constraint-mediated routing model in which
what gets transported where is determined by the topology of admissibility constraints
at each stage of the transport system.

The appropriate intermediate framework for making this connection precise is that of
graph routing under local compatibility constraints. Model the ER exit site network as
a directed graph $G = (V, E)$ in which nodes represent distinct membrane domains
(rough ER, smooth ER, transitional zones, exit sites) and directed edges represent
possible vesicular transfer events. Each edge $e = (u, v)$ carries a compatibility
predicate $\chi_e : \mathcal{V} \to \{0,1\}$ specifying which cargo configurations
are permissible on that route. The routing problem is then a constraint satisfaction
problem on $G$: find a flow $f : E \to \Rspace_{\geq 0}$ that satisfies capacity
constraints, flow conservation, and the compatibility predicates $\chi_e$ simultaneously.
This is precisely the structure of queueing geometry under compatibility constraints:
the ER export machinery is not a pipeline but a constraint satisfaction engine, and
the ER exit site ultrastructure encodes the compatibility constraints that make
selective, reliable transport possible. The RSVP vector field $\vfield$ on the cellular
domain then corresponds to the induced flow field of the routing solution, and the
admissibility manifold of the transport system is the polytope of feasible flows
satisfying all compatibility constraints.


\begin{definition}[Intracellular Transport Admissibility]
Let $\mathcal{V}$ be the space of possible vesicular cargo configurations and
$\mathcal{R} \subset \mathcal{V}$ the set of cargo configurations admissible for
a given exit site geometry (satisfying coat protein binding compatibility, size
constraints, and lipid composition requirements). The transport routing function
is the map $\tau : \mathcal{V} \to 2^{\text{Destinations}}$ defined by
$\tau(v) = \{d : v \in \mathcal{R}_d\}$, where $\mathcal{R}_d$ is the admissibility
set for destination $d$.
\end{definition}

\subsection{Ribosome Exclusion and Spatial Constraint}

A key structural finding of~\cite{Downes2026ERES} is that ER exit sites maintain
ribosome exclusion zones: regions from which ribosomes are geometrically excluded,
enabling membrane deformation and vesicle budding without interference from the
translation machinery. This is a physical realization of the RSVP vorticity constraint:
the exit site geometry imposes a spatial constraint on the vector field of intracellular
flow, preventing turbulent interference between translation and export.

\begin{proposition}[Exclusion Zones as Admissibility Enforcement]
The ribosome exclusion zone around an ER exit site of radius $r$ enforces the local
admissibility condition
\begin{equation}
  \norm{\grad \times \vfield_\text{ribosomes}}_{B(x_0, r)} < \varepsilon
\end{equation}
for the ribosomal flow field $\vfield_\text{ribosomes}$ in a ball of radius $r$
around the exit site center $x_0$, ensuring that membrane deformation dynamics are
not disrupted by competing mechanical forces.
\end{proposition}

\subsection{Semantic Infrastructure and Cellular Logistics}

The transport topology revealed by~\cite{Downes2026ERES} is a biological instance of
what we call semantic infrastructure: a logistics system in which the routing of
information-bearing entities is determined not by explicit addressing but by
compatibility constraints encoded in the geometric structure of the transport medium
itself. The cell does not instruct vesicles where to go; it constructs a topological
manifold of admissible routes, and vesicle destinations emerge from the intersection
of cargo properties with route geometry.

This provides a molecular precedent for distributed semantic routing in artificial
systems. An inference architecture modeled on ER exit site topology would route
information-bearing representations through a network of compatibility-constrained
junctions rather than through explicit addressable channels.

%% ====================================================================
\section{Neuromorphic Matter: Morphology as Computation}
%% ====================================================================

\subsection{Stretchable Neuromorphic Circuits}

The large-scale stretchable neuromorphic circuit developed by Li and
collaborators~\cite{Li2026Stretchable} represents a physical realization of the
hypothesis that cognition need not be localized in a rigid processor but can be
distributed across a deformable material substrate. Their system integrates synaptic
transistors, artificial neurons, and spike-based communication pathways into a
stretchable polymer matrix, enabling on-body edge computing that deforms continuously
with the body it monitors.

The conceptual significance of this work extends far beyond its engineering achievements.
It demonstrates that the substrate of computation can be morphological: the geometry
and deformability of the material medium are active components of inference, not
passive supports. When the body bends, the neuromorphic circuit bends with it, and the
computation adapts to the changed geometry without explicit reconfiguration. Morphology
is computation.

\begin{definition}[Xylomorphic Computation]
A computational system is xylomorphic if its computational dynamics are functions of
the geometry of its material substrate, formally: if the effective Hamiltonian
$H_\text{eff}$ governing the system's computation satisfies
\begin{equation}
  H_\text{eff} = H_0 + \int_\Omega h[\mathbf{u}(\mathbf{x})] \,d\mathbf{x},
\end{equation}
where $\mathbf{u}(\mathbf{x})$ is the displacement field of the substrate and
$h[\mathbf{u}]$ encodes the morphological coupling between substrate geometry and
computational dynamics.
\end{definition}

\subsection{Distributed Constraint Fields in Material Substrates}

The RSVP interpretation of the stretchable neuromorphic system treats the material
substrate as a physical instantiation of the RSVP field triple. The substrate's elastic
properties correspond to the density field $\Phi$, encoding the local stiffness and
energy density of the computation medium. The ionic or electronic signal pathways
correspond to the vector flow field $\vfield$, encoding directed information transport.
The thermal and entropic properties of the semiconductor interfaces correspond to the
entropy field $S$, encoding noise and dissipation characteristics.

Under this interpretation, the neuromorphic computation performed by the stretchable
circuit is precisely a constrained trajectory in the RSVP admissibility manifold: the
system's computation must remain within the admissibility set defined by its material
properties (mechanical integrity, signal-to-noise bounds, power dissipation limits).

\begin{theorem}[\ClaimH\ RSVP Reduction of Xylomorphic Computation]
\label{thm:xylomorphic}
Let $(\Phi, \vfield, S)$ be an RSVP field triple on a substrate manifold $\Omega$
with deformation field $\mathbf{u}$. If the computation performed by the system is
represented as a trajectory $\gamma : [0,T] \to \Adm$, then the computational capacity
of the system, measured by the Shannon entropy of the output distribution, satisfies
\begin{equation}
  C_\text{comp} \leq \log |\pi_0(\Adm_\mathbf{u})|
\end{equation}
where $\Adm_\mathbf{u}$ is the admissibility manifold deformed by $\mathbf{u}$ and
$\pi_0$ denotes the set of path-connected components.
\end{theorem}

\begin{proof}
Each path-connected component of $\Adm_\mathbf{u}$ supports a distinct class of
system trajectories. Trajectories in different components are separated by
inadmissibility barriers and therefore produce distinguishable outputs. The
Shannon entropy of the output distribution is bounded by the logarithm of the number
of distinguishable trajectory classes, which equals $|\pi_0(\Adm_\mathbf{u})|$.
\end{proof}

Theorem~\ref{thm:xylomorphic} has a concrete implication: the computational capacity
of a xylomorphic system is a topological invariant of the deformed admissibility
manifold. Deforming the substrate changes the topology of $\Adm_\mathbf{u}$ and
thereby changes the system's computational capacity. Morphology is not merely a physical
convenience but a topological determinant of what the system can compute.

%% ====================================================================
\section{Nonlinear Optomechanics: Coherence at Threshold}
\label{sec:opto}
%% ====================================================================

\subsection{Self-Sustained Oscillation in the Low-Excitation Regime}

The nonlinear optomechanical system studied by Dhiman and collaborators~\cite{Dhiman2026Optomechanics}
demonstrates self-sustained oscillation at ultra-low excitation levels through
Kerr-enhanced superconducting coupling. Their system consists of a mechanical resonator
coupled to a superconducting microwave cavity with Kerr nonlinearity, and achieves
sustained oscillation far below the threshold that would be required in a linear system.
The Kerr nonlinearity effectively lowers the admissibility threshold, enabling coherent
behavior at energy scales where a linear system would simply thermalize.

This is a direct physical realization of the RSVP admissibility threshold concept.
The Kerr nonlinearity modifies the effective admissibility manifold $\Adm$ by introducing
a nonlinear correction to the pressure functional $P(\Phi, S)$ in equation~\eqref{eq:rsvp2},
expanding the basin of attraction for coherent oscillatory trajectories.

\begin{definition}[Kerr-Enhanced Admissibility]
In a Kerr-coupled optomechanical system, the effective pressure functional is modified as
\begin{equation}
  P_\text{Kerr}(\Phi, S) = P_0(\Phi, S) - \chi_K \Phi^2,
\end{equation}
where $\chi_K > 0$ is the Kerr coupling constant. The Kerr-enhanced admissibility
manifold $\Adm_\text{Kerr}$ is the admissibility manifold defined using $P_\text{Kerr}$
in place of $P_0$.
\end{definition}

\begin{proposition}[\ClaimR\ Threshold Reduction by Kerr Nonlinearity]
The oscillation onset threshold in $\Adm_\text{Kerr}$ is strictly lower than in
$\Adm_0$ (the linear system), satisfying
\begin{equation}
  \Phi_\text{threshold}^\text{Kerr} = \Phi_\text{threshold}^0 - \chi_K \cdot \Delta\Phi
\end{equation}
for a geometry-dependent constant $\Delta\Phi > 0$, confirming the experimental
observation of sub-threshold oscillation in~\cite{Dhiman2026Optomechanics}.
\end{proposition}

\subsection{Recursive Stabilization and Attractor Formation}

The self-sustained oscillation described in~\cite{Dhiman2026Optomechanics} is an
example of a limit cycle attractor: a periodic orbit in the phase space of the
optomechanical system toward which nearby trajectories converge. In the RSVP framework,
this attractor corresponds to a periodic trajectory confined to the admissibility manifold.

The TARTAN framework provides a tiling interpretation of this attractor structure.
The phase space of the optomechanical system is tiled by trajectory patches, each
annotated with noise statistics. The limit cycle is the unique trajectory that closes
under the recursive tiling operation: each segment of the cycle tiles forward to the
next segment, and the cycle closes when the tiling returns to its starting configuration.

\begin{definition}[TARTAN Cycle]
A TARTAN cycle is a closed trajectory $\gamma : [0,T] \to \Adm$ with $\gamma(T) = \gamma(0)$
such that the trajectory tiling $\mathcal{T}[\gamma]$ is periodic: there exists a minimal
period $T_0 | T$ such that $\mathcal{T}[\gamma](t) = \mathcal{T}[\gamma](t + T_0)$ for
all $t$.
\end{definition}

%% ====================================================================
\section{Morphology as an Inference Prior}
\label{sec:morphprior}
%% ====================================================================

Before turning to specific applications of structured priors in graph neural networks,
clinical AI, and ontology-enhanced classification, it is worth articulating the
general principle unifying those three cases. Each one represents a setting in which
the morphology of the domain---its relational geometry, its diagnostic topology, its
semantic organization---functions not merely as background context for inference but
as a \emph{prior} that regularizes the hypothesis space and eliminates inadmissible
conclusions before any statistical evidence is considered.

The Bayesian framing makes this precise. Let $\mathcal{H}$ be the hypothesis space
and $p_0$ the prior distribution over $\mathcal{H}$. In an unconstrained statistical
system, $p_0$ is uniform or weakly regularized (e.g., by $L^2$ penalties). The
posterior $p(h | \text{data})$ is therefore determined almost entirely by the
likelihood term. In a morphologically constrained system, $p_0$ is instead concentrated
on the admissibility submanifold $\mathcal{H}_\Adm \subset \mathcal{H}$:
\begin{equation}
  p_0(h) \propto \mathbf{1}[h \in \mathcal{H}_\Adm] \cdot \exp(-\beta \cdot
  \text{dist}(h, \partial \mathcal{H}_\Adm)).
\end{equation}
The admissibility prior is not a soft preference but a hard filter: hypotheses
outside $\mathcal{H}_\Adm$ receive negligible prior weight regardless of their
likelihood. This hard filter substantially reduces the effective dimension of the
hypothesis space and concentrates posterior mass on the morphologically coherent region.

The performance improvements documented across all three subsequent sections ---
pathological topology priors in whole slide image classification, symbolic dermatological
criteria in BCC diagnosis, and ontological features in Amharic news classification ---
are all instances of this mechanism. In each case, incorporating morphological structure
into the prior reduces effective hypothesis space dimension, eliminates inferential
noise from inadmissible regions, and concentrates posterior mass on the semantically
or diagnostically meaningful submanifold.

\begin{proposition}[\ClaimR\ Morphological Prior Reduces Effective Dimension]
\label{prop:morphdim}
Let $\mathcal{H}$ be a hypothesis space of dimension $n$ and $\mathcal{H}_\Adm
\subset \mathcal{H}$ an admissible submanifold of dimension $k < n$. Under an
admissibility prior $p_0$ concentrated on $\mathcal{H}_\Adm$, the effective number
of parameters governing inference satisfies
\begin{equation}
  n_\text{eff} \leq k + O\!\left(\frac{1}{\beta}\right),
\end{equation}
where $\beta$ is the concentration parameter of $p_0$, and the classification risk
satisfies $R_\Adm \leq R_\mathcal{H} \cdot (k/n) + O(1/N)$ for $N$ training samples.
\end{proposition}

\begin{proof}
The effective dimension $n_\text{eff}$ is the trace of the Fisher information matrix
restricted to the support of $p_0$, which is concentrated on $\mathcal{H}_\Adm$.
The trace of the Fisher matrix on a $k$-dimensional submanifold equals $k$ in the
limit $\beta \to \infty$, with corrections of order $1/\beta$ from the soft boundary
of $p_0$. The risk bound follows from the bias-variance decomposition with effective
dimension $k$: variance scales as $k/N$, and bias is zero on $\mathcal{H}_\Adm$
(since the ground truth is admissible by assumption).
\end{proof}

Proposition~\ref{prop:morphdim} provides a precise information-theoretic account of
why morphological priors improve inference. They do not add new statistical evidence;
they reduce the search space. The improvement scales with the ratio $k/n$: the more
informative the morphological prior (i.e., the lower $k$ relative to $n$), the
greater the performance gain. This is precisely the mechanism underlying all three
sections that follow.

%% ====================================================================
\section{Graph Neural Networks and Pathological Topology Priors}
\label{sec:gnn}
%% ====================================================================

\subsection{The Failure of Pure Statistical Aggregation}

The pathology-prior driven graph neural network of Wu and Jiang~\cite{Wu2026GNN} for
whole slide image (WSI) classification confronts a fundamental limitation of standard
graph neural network architectures: when applied to histopathological tissue without
explicit structural guidance, they fail to leverage diagnostically meaningful tissue
configurations and instead aggregate local features without regard for the relational
geometry that pathologists use to reason about disease.

This failure is the empirical instantiation of a theoretical claim central to the RSVP
and CLIO frameworks: statistical aggregation over unconstrained feature distributions
cannot recover the structural distinctions encoded in admissibility geometry. What
matters diagnostically in pathological tissue is not the distribution of individual
cellular features but the topology of relational configurations among cells, glands,
and tissue compartments.

\begin{definition}[Pathological Admissibility Prior]
Let $G = (V, E)$ be a tissue graph with node set $V$ (cells or regions) and edge set
$E$ (spatial adjacency or functional interaction). A pathological admissibility prior
is a constraint set $\mathcal{P} \subset 2^E$ specifying which edge subgraphs
correspond to diagnostically meaningful tissue configurations (e.g., glandular formations,
invasion fronts, stromal patterns). A message-passing operation $\mu : \mathcal{F}^{|V|}
\to \mathcal{F}^{|V|}$ is $\mathcal{P}$-admissible if it aggregates features only along
edges $e \in E'$ for some $E' \in \mathcal{P}$.
\end{definition}

\subsection{CLIO Formalization of Topology-Aware GNNs}

The CLIO framework provides a natural formalization of the pathology-prior approach.
The classification task is a constraint satisfaction problem: find a representation
$h : G \to \mathcal{Y}$ from tissue graphs to diagnostic labels such that $h$ is
consistent with the pathological admissibility prior $\mathcal{P}$ and maximizes
diagnostic accuracy.

\begin{theorem}[\ClaimR\ CLIO Optimality of Topology-Prior GNNs]
\label{thm:gnn-clio}
Let $\mathcal{H}$ be the class of all graph neural networks and $\mathcal{H}_\mathcal{P}
\subset \mathcal{H}$ the subclass of $\mathcal{P}$-admissible graph neural networks.
For any diagnostic task with true label distribution $p(y|G)$ that depends on
$\mathcal{P}$-admissible substructures, the optimal risk in $\mathcal{H}_\mathcal{P}$
is strictly lower than the optimal risk in $\mathcal{H} \setminus \mathcal{H}_\mathcal{P}$.
\end{theorem}

\begin{proof}
The true label distribution $p(y|G)$ depends on $\mathcal{P}$-admissible substructures
by hypothesis, meaning there exists a function $f : 2^E \to \mathcal{Y}$ such that
$y = f(E')$ for some $E' \in \mathcal{P}$ with probability 1 under $p$. Any
$\mathcal{P}$-admissible GNN can represent such a function exactly (as its message
passing is restricted to $\mathcal{P}$-admissible subgraphs), whereas any GNN in
$\mathcal{H} \setminus \mathcal{H}_\mathcal{P}$ aggregates over non-admissible subgraphs,
introducing irreducible noise from pathologically irrelevant edge configurations.
The risk difference is bounded below by the mutual information $I(y; E_\text{non-adm})$
between the true label and non-admissible edge features, which is strictly positive
by the dependence assumption.
\end{proof}

\subsection{Substructure Awareness as Constraint Sensitivity}

The practical implication of Theorem~\ref{thm:gnn-clio} is that incorporating
pathological priors into GNN architecture is not merely a heuristic improvement
but a theoretically necessary step for achieving optimal diagnostic accuracy on
tasks where the ground truth depends on admissible relational structure. The fact
that Wu and Jiang~\cite{Wu2026GNN} demonstrate substantial empirical improvements
from their topology-aware architecture is precisely what the CLIO framework predicts.

%% ====================================================================
\section{Clinical AI and Symbolic-Statistical Integration}
%% ====================================================================

\subsection{Dual Explanation in BCC Diagnosis}

The multitask learning system for basal cell carcinoma (BCC) diagnosis by Matas and
collaborators~\cite{Matas2026BCC} demonstrates that clinically useful artificial
intelligence requires encoding diagnostically meaningful symbolic structure into
network architecture rather than treating explanation as a post hoc interpretation
of opaque statistical decisions. Their dual explanation architecture simultaneously
generates a dermoscopic label explanation and a histopathological structure explanation,
both grounded in the symbolic vocabulary of clinical dermatology.

This is an architectural realization of the RSVP admissibility principle at the
epistemological level. The network's predictions are constrained to lie within the
admissibility manifold of clinically meaningful explanations: a prediction is admissible
only if it is consistent with recognized dermoscopic criteria and histopathological
patterns. The dual explanation mechanism ensures that the network's inference trajectory
remains within this admissibility manifold throughout the classification process.

\begin{definition}[Clinically Admissible Explanation]
Let $\mathcal{E}_D$ be the space of dermoscopic explanations and $\mathcal{E}_H$ the
space of histopathological explanations, each grounded in established clinical
ontologies. An explanation pair $(e_D, e_H) \in \mathcal{E}_D \times \mathcal{E}_H$ is
clinically admissible if the two explanations are mutually consistent (the dermoscopic
features are causally compatible with the histopathological findings) and individually
grounded (each explanation uses only clinically recognized feature categories).
\end{definition}

\subsection{Simulated Agency and Intelligible Inference}

The dual explanation architecture of~\cite{Matas2026BCC} instantiates what the
Simulated Agency framework calls intelligible inference: inference whose process, not
merely whose output, can be mapped onto a symbolic structure recognizable to an expert
agent. This requires that the network's internal dynamics project onto the clinically
admissible explanation manifold at each inference step, not merely at the final
output layer.

\begin{theorem}[\ClaimH\ Explanation Consistency as Trajectory Constraint]
\label{thm:explanation}
Let $f : \mathcal{X} \to \mathcal{Y}$ be a classifier with intermediate representation
$z \in \mathcal{Z}$. The classifier produces clinically admissible explanations for
all inputs if and only if the inference trajectory $x \mapsto z \mapsto y$ is confined
to the admissibility manifold $\Adm_\mathcal{Z} \subset \mathcal{Z}$ defined by the
joint clinical admissibility constraint.
\end{theorem}

This theorem establishes explanation as a trajectory constraint rather than a
post hoc decoration. An explanatory architecture is one whose inference dynamics
are geometrically constrained to the space of admissible explanations, not one that
generates post hoc justifications for unconstrained inference.

%% ====================================================================
\section{Ontology-Constrained Inference and Semantic Geometry}
%% ====================================================================

\subsection{Amharic News Classification and Structural Priors}

The ontology-enhanced Amharic news classification system of Taye and
collaborators~\cite{Taye2026Amharic} provides an unusually clean empirical
demonstration of the constraint-before-content principle at the semantic level.
Their finding is that incorporating structured ontological features into a TF-IDF
baseline classifier substantially improves classification accuracy over the pure
statistical baseline. The ontology encodes semantic relationships among concepts,
providing a symbolic structure that constrains the admissible classification trajectories.

The result is striking because the baseline system already has access to all the
statistical information available in the text. The ontological augmentation does not
add new textual evidence; it adds structural constraints on how that evidence can be
combined. The performance improvement is entirely attributable to the imposition of
admissibility constraints on the inference process.

\begin{definition}[Semantic Admissibility via Ontology]
Let $\mathcal{O} = (C, R, A)$ be an ontology with concept set $C$, relation set $R$,
and axiom set $A$. For a text representation $\mathbf{x} \in \Rspace^d$, the
ontology-admissible classification set is
\begin{equation}
  \mathcal{Y}_\mathcal{O}(\mathbf{x}) = \left\{ y \in \mathcal{Y} : \exists c \in C \text{ s.t. }
  c \vdash y \text{ under } A \text{ and } \text{sim}(\mathbf{x}, \mathbf{e}_c) > \theta \right\},
\end{equation}
where $\mathbf{e}_c$ is the feature embedding of concept $c$, $\text{sim}$ is a
semantic similarity measure, and $\theta$ is a confidence threshold.
\end{definition}

\subsection{RSVP Interpretation of Semantic Topology}

In the RSVP framework, an ontology $\mathcal{O}$ defines a topological structure on
the semantic manifold. The concepts $C$ are nodes, the relations $R$ are edges, and
the axioms $A$ specify the admissibility constraints on semantic trajectories. A
classification system that respects the ontological structure is one whose inference
trajectories remain within the semantic admissibility manifold defined by $\mathcal{O}$.

\begin{proposition}[\ClaimP\ Ontological Constraints Reduce Classification Error]
Let $p_0$ be the classification error of a purely statistical classifier and $p_\mathcal{O}$
the classification error of the ontology-augmented classifier. Under the assumption that
the true label distribution is consistent with $\mathcal{O}$ (i.e., the ground-truth
labels respect the ontological axioms $A$), we have
\begin{equation}
  p_\mathcal{O} \leq p_0 - \delta,
\end{equation}
where $\delta > 0$ is determined by the mutual information between the ontological
structure and the label distribution.
\end{proposition}

This proposition formalizes the empirical finding of~\cite{Taye2026Amharic}: the
performance improvement from ontological augmentation is not accidental but reflects
the systematic elimination of classification hypotheses inconsistent with the
admissibility structure of the task's semantic domain.

%% ====================================================================
\section{Topology Versus Localization}
\label{sec:topovslocal}
%% ====================================================================

The tremor network paper of the following section is one of several in this corpus
that implicitly argue against a deep assumption embedded in much of twentieth-century
systems biology and neuroscience: the assumption that complex functional properties
are \emph{localized} in specific components. On this view, a symptom is a property
of a damaged region; a function is a property of a dedicated module; a failure is a
property of a single faulty node.

The localizationist framework is attractive because it licenses targeted intervention:
if a symptom lives at a location, removing or stimulating that location should
address the symptom. But the evidence reviewed in this paper suggests that
localizationism is systematically wrong as an account of complex system behavior,
and the RSVP framework provides a formal account of why.

In an RSVP field system, the field triple $(\Phi, \vfield, S)$ is defined globally on
$M$. No property of the system is localized at a single point: every quantity depends
on the global field configuration, and perturbations at one location propagate
nonlocally through the coupled field equations. A symptom---an anomalous pattern in
$\Phi$, a disrupted flow pattern in $\vfield$, an entropy spike in $S$---is a property
of a field configuration, not of a point.

The contrast with localizationist ontology is sharp. A localizationist account of
a symptom $\sigma$ identifies a point $x_0 \in M$ such that $\sigma$ is caused by
an anomaly at $x_0$. The RSVP account identifies a region $U_\sigma \subset
\mathcal{X}$ in field configuration space---the distributed causal manifold $\mathcal{M}_\sigma$
from Definition~11.1 below---such that field configurations in $U_\sigma$ produce $\sigma$.
The two accounts make different predictions about intervention: the localizationist
predicts that targeting $x_0$ resolves $\sigma$; the RSVP account predicts that
the effective intervention is any perturbation of the field that exits $U_\sigma$.
The empirical data on DBS strongly favor the second account.

More generally, the topology-versus-localization distinction applies across every
domain in this paper. Protein nanocage geometry is a property of the closure manifold,
not of individual monomers. ER exit site routing is a property of the transport graph
topology, not of individual vesicles. Neural coherence is a property of the spectral
geometry of the weight matrix, not of individual neurons. Diagnostic accuracy in
pathological AI is a property of the topology-aware message-passing architecture,
not of individual feature detectors.

In each case, the localizationist account fails to capture the relevant organizational
level because organization is a topological rather than a pointwise property. This
is not a limitation of current scientific instruments or resolution; it is a
fundamental feature of constraint-mediated systems. The organization is in the
manifold, not in the matter.

%% ====================================================================
\section{Tremor Networks and Distributed Causal Manifolds}
\label{sec:tremor}
%% ====================================================================

\subsection{Symptom Expression as Network Topology}

The tremor and DBS network study by Weigl and collaborators~\cite{Weigl2026Tremor}
demonstrates that tremor expression and its suppression by deep brain stimulation are
properties of distributed metabolic and functional network topology rather than
localized pathological defects. Their analysis shows that converging metabolic and
functional networks determine both the expression of tremor and the therapeutic
efficacy of DBS, with no single brain region uniquely responsible for either.

This finding is a direct empirical refutation of localizationist causality and a
confirmation of the distributed causal manifold model implicit in the RSVP framework.
Symptoms are not properties of individual brain regions but of network configurations:
they emerge from the topology of the admissibility manifold of the functional brain
network, and therapy acts by modifying that topology rather than by targeting a
localized defect.

\begin{definition}[Distributed Causal Manifold]
The distributed causal manifold of a symptom $\sigma$ is the set
\begin{equation}
  \mathcal{M}_\sigma = \left\{ (F, M) \in \mathcal{F} \times \mathcal{M} :
  \Pr[\sigma | F, M] > \theta_\sigma \right\},
\end{equation}
where $\mathcal{F}$ is the space of functional connectivity matrices, $\mathcal{M}$ is
the space of metabolic network states, and $\theta_\sigma$ is a symptom expression
threshold.
\end{definition}

\subsection{DBS as Topological Modification}

Under the distributed causal manifold model, deep brain stimulation acts not by
suppressing a localized pathological signal but by modifying the topology of the
functional network so that the network state exits $\mathcal{M}_\sigma$. The therapeutic
target is not a brain region but a network configuration, and the therapeutic mechanism
is not signal suppression but topological redirection.

\begin{theorem}[\ClaimH\ DBS as Admissibility Redirection]
\label{thm:dbs}
Let $\mathcal{M}_\sigma$ be the distributed causal manifold of tremor. Effective DBS
stimulation is characterized by a stimulation pattern that maps the current network
state $x \in \mathcal{M}_\sigma$ to a state $x' \notin \mathcal{M}_\sigma$ while
minimizing the total displacement $\norm{x' - x}$ in functional network space.
\end{theorem}

\begin{proof}
If DBS acts by direct suppression of a localized signal, it would need to eliminate
the signal from a single node, which would generically fail due to network redundancy
(other nodes would compensate). The empirical success of DBS at distributed stimulation
targets therefore implies that it acts by network-level modification, i.e., by displacing
the network state in functional space. The optimal displacement is the minimum-norm
displacement that exits $\mathcal{M}_\sigma$, which corresponds to moving perpendicular
to the boundary $\partial \mathcal{M}_\sigma$.
\end{proof}

The prediction of Theorem~\ref{thm:dbs}---that optimal DBS targets are determined by
the geometry of the boundary $\partial \mathcal{M}_\sigma$ rather than by anatomical
proximity to a lesion---is consistent with the clinical observation that DBS efficacy
depends on stimulation of specific network nodes rather than proximity to dopaminergic
pathways.

%% ====================================================================
\section{Paleoclimate Reconstruction: Memory as Distributed Trajectory Residue}
%% ====================================================================

\subsection{Rhodolith Archives and Fragmented Temporal Records}

The rhodolith paleotemperature study by Li and collaborators~\cite{Li2026Rhodoliths}
demonstrates that daily-resolution temperature records extending back thousands of years
can be reconstructed from the growth banding patterns of coralline algae. Their key
methodological contribution is the use of dynamic time warping (DTW) to align and
combine partial growth records from multiple rhodolith specimens, reconstructing a
coherent temporal history from fragmentary evidence.

This methodology instantiates a model of memory that is deeply compatible with the
RSVP and TARTAN frameworks: memory as distributed structural residue across overlapping
trajectories, rather than as discrete stored states. The complete temperature history
does not exist in any single rhodolith; it exists in the statistical overlap structure
of many partial records, and can be recovered by a constraint-sensitive alignment process.

\begin{definition}[Distributed Temporal Memory]
Let $\{\gamma_i : [t_i^-, t_i^+] \to \mathcal{X}\}_{i=1}^N$ be a collection of
partial trajectories in a state space $\mathcal{X}$, with potentially overlapping
temporal domains. The distributed temporal memory $\mathcal{M}_T$ is the unique
trajectory $\hat\gamma : [t_\text{min}, t_\text{max}] \to \mathcal{X}$ defined by
\begin{equation}
  \hat\gamma(t) = \arg\min_{x \in \mathcal{X}} \sum_{i : t \in [t_i^-, t_i^+]}
  d(\gamma_i(t), x)^2,
\end{equation}
where $d$ is the metric on $\mathcal{X}$ and the sum runs over all partial trajectories
whose temporal domain includes $t$.
\end{definition}

\subsection{Dynamic Time Warping as Constraint Closure}

The dynamic time warping algorithm used in~\cite{Li2026Rhodoliths} solves a discrete
version of the distributed temporal memory problem: it finds the alignment between two
partial sequences that minimizes the total distance between aligned elements, subject
to the monotonicity constraint that the alignment preserves temporal order.

In the TARTAN framework, this alignment operation corresponds to constraint closure
over the tiling of the temporal state space. Each rhodolith growth sequence is a partial
tiling of the paleoclimate trajectory, and DTW finds the unique closed tiling consistent
with all partial tilings simultaneously.

\begin{proposition}[\ClaimP\ DTW as TARTAN Constraint Closure]
The dynamic time warping alignment of $N$ partial trajectory records $\{\gamma_i\}$
is the solution to the TARTAN constraint closure problem
\begin{equation}
  \hat\gamma = \arg\min_{\gamma \in \mathcal{T}[\mathcal{X}]}
  \sum_{i=1}^N \text{DTW}(\gamma|_{[t_i^-, t_i^+]}, \gamma_i),
\end{equation}
where $\mathcal{T}[\mathcal{X}]$ denotes the space of TARTAN-admissible trajectories
in $\mathcal{X}$ (trajectories consistent with the smoothness and range constraints
of the paleoclimate system).
\end{proposition}

\subsection{Epistemological Implications}

The distributed temporal memory model has broader implications for the epistemology of
historical reconstruction. Any system whose history must be reconstructed from partial
and overlapping records---geological systems, biological lineages, computational logs,
institutional histories---faces the same structural problem: extracting a coherent
trajectory from fragmentary evidence. The TARTAN framework provides a general solution:
treat each partial record as a constraint on the set of admissible global trajectories,
and identify the unique trajectory satisfying all constraints simultaneously.

%% ====================================================================
\section{Universality Across Domains}
\label{sec:universality}
%% ====================================================================

The preceding ten sections have surveyed a wide range of empirical systems, all
interpreted through the lens of admissibility geometry. Before synthesizing these
results into a single formal principle, it is worth addressing a potential objection:
is the recurrence of admissibility-structured behavior across such different systems
genuinely a mathematical discovery, or is it an artifact of a sufficiently flexible
interpretive framework?

The objection has force. The RSVP field triple $(\Phi, \vfield, S)$ and the
admissibility manifold $\Adm$ are general enough to be instantiated by almost any
constrained dynamical system. One might worry that the cross-domain correspondences
documented in this paper are tautological: of course a framework flexible enough to
model neural fields, protein folding, and paleoclimate reconstruction will find its
own structure in each.

The response to this objection is threefold. First, the correspondences are not
merely terminological but predictive. The identification of neural criticality with
vanishing spectral abscissa (Theorem~\ref{thm:clio-criticality}) makes a specific
quantitative prediction---that systems tuned to spectral radius 1 will exhibit
power-law eigenspectra---that was subsequently confirmed by Pachitariu et al.~\cite{Pachitariu2026Critical}
independently of the RSVP framework. The identification of protein nanocage geometry
with constrained configuration spaces (Theorem~\ref{thm:nanocage}) predicts specific
dimension counts for the accessible symmetry classes, which can be tested against the
combinatorial geometry of polyhedral assemblies. These predictions are not tautological.

Second, the framework exhibits non-trivial universality in the renormalization group
sense. Different systems---neural networks, protein assemblies, optomechanical
resonators---are governed by microscopically different equations but share the same
universality class of behavior near their admissibility boundaries. This universality
is the hallmark of a genuine structural invariant, not merely a flexible vocabulary.
Two systems belong to the same universality class if and only if they share the same
symmetry group acting on their admissibility manifold. The correspondence, therefore,
is a statement about symmetry, not a statement about fitting parameters.

Third, and most importantly, the framework makes a \emph{negative} prediction that
distinguishes it from vacuous universalism: systems that lack admissibility structure
should \emph{not} exhibit the organizational phenomena documented in this paper.
Purely random networks without constraint structure should not exhibit long-range
coherence. Purely statistical classifiers without structural priors should perform
worse than topology-aware classifiers. This negative prediction is falsifiable and
is supported by the empirical contrasts reported in~\cite{Wu2026GNN},~\cite{Taye2026Amharic},
and~\cite{Matas2026BCC}, where unconstrained baselines consistently underperform
admissibility-constrained architectures.

The universality across domains, then, is not an artifact of interpretive flexibility.
It reflects the mathematical fact that any system maximizing long-range correlation
subject to stability constraints will be drawn toward the boundary of its admissibility
manifold, and any system near such a boundary will exhibit the universal scaling
behavior described by Proposition~\ref{prop:powerlaw} and
Conjecture~\ref{conj:criticality}. The domains are different; the geometry is the same.

%% ====================================================================
\section{The Admissibility Principle: A Unified Account}
%% ====================================================================

\subsection{Statement of the Principle}

The ten papers reviewed in this work, spanning neural dynamics, molecular biology,
cellular transport, materials science, clinical AI, computational linguistics, and
paleoclimate reconstruction, converge on a single structural finding that can be
stated as a general principle.

\begin{tcolorbox}[rsvpbox, title={The Admissibility Principle}]
\textbf{Admissibility Principle.} In any system capable of sustaining coherent
global organization over extended time scales, the geometry of admissible configurations
and admissible transitions is the primary organizational datum. Content---specific
states, signals, or molecular identities---is secondary: it is constrained by and
derivative of admissibility geometry. Global coherence emerges not from explicit
centralized control but from the local dynamics of a system confined to its
admissibility manifold.
\end{tcolorbox}

\subsection{Mathematical Formalization}

We now state and prove the main theorem of the paper, which formalizes the admissibility
principle as a precise mathematical claim.

\begin{theorem}[\ClaimH\ Admissibility Determines Coherence Class]
\label{thm:main}
Let $\mathcal{S}$ be a dynamical system on a state space $\mathcal{X}$ with flow
$\phi_t : \mathcal{X} \to \mathcal{X}$, and let $\Adm \subset \mathcal{X}$ be an
admissibility manifold invariant under $\phi_t$. Suppose two systems $\mathcal{S}_1$
and $\mathcal{S}_2$ have distinct local dynamics but identical admissibility manifolds
$\Adm_1 = \Adm_2 = \Adm$. Then:
\begin{enumerate}[label=(\roman*)]
  \item The long-run statistical behavior of $\mathcal{S}_1$ and $\mathcal{S}_2$ are
    equivalent: $\mu_1 = \mu_2$ where $\mu_i$ is the ergodic measure of $\phi_t^{(i)}$
    restricted to $\Adm$.
  \item The global coherence properties of $\mathcal{S}_1$ and $\mathcal{S}_2$---defined
    as the scaling exponents of spatial and temporal correlation functions---are identical.
  \item Any observable $f : \mathcal{X} \to \Rspace$ that distinguishes $\mathcal{S}_1$
    from $\mathcal{S}_2$ must be sensitive to the transient dynamics within $\Adm$ rather
    than to its geometry.
\end{enumerate}
\end{theorem}

\begin{proof}
(i) Since $\Adm$ is invariant and both systems are confined to $\Adm$, the ergodic
measures $\mu_i$ are supported on $\Adm$. By the uniqueness of the ergodic measure for
a uniquely ergodic system on a compact manifold, $\mu_1 = \mu_2$ when $\Adm_1 = \Adm_2$.
For non-uniquely ergodic systems, the ergodic decomposition is identical since it is
determined by the topology of $\Adm$.

(ii) Coherence scaling exponents are determined by the spectrum of the transfer operator
of the flow restricted to $\Adm$, which depends only on the geometry of $\Adm$ and not
on the specific flow dynamics within it (by the universality of critical exponents in
systems with the same symmetry class).

(iii) If $f$ distinguishes $\mathcal{S}_1$ from $\mathcal{S}_2$ at the level of their
stationary distributions, then $\int f d\mu_1 \neq \int f d\mu_2$. But (i) shows
$\mu_1 = \mu_2$, so no such $f$ exists. Therefore any distinguishing observable must
depend on the transient behavior.
\end{proof}

\begin{remark}[Scope of Claim~(i)]
Statement~(i) invokes uniqueness of the ergodic measure, which holds for uniquely
ergodic systems (e.g.\ minimal flows on compact manifolds) but fails in general for
systems with multiple invariant measures. For the cross-domain applications in
Corollary~\ref{cor:domains}, the admissibility manifolds in question are compact and
the confined dynamics are generically uniquely ergodic by the structural stability
theorems of Mañé and Bowen. For systems where unique ergodicity cannot be assumed,
statement~(i) should be read in the weaker sense: the \emph{ergodic decompositions}
of $\mu_1$ and $\mu_2$ are identical, meaning the two systems have the same family
of ergodic invariant measures supported on $\Adm$, even if neither measure is unique.
Statements~(ii) and~(iii) are unaffected by this qualification.
\end{remark}

\subsection{Cross-Domain Instances of the Admissibility Principle}

Theorem~\ref{thm:main} has a specific implication: systems with the same admissibility
geometry will exhibit the same global coherence properties, regardless of their
microscopic details. The following corollary collects the instances demonstrated
across the ten papers reviewed.

\begin{corollary}[\ClaimC\ Cross-Domain Admissibility Equivalence]
\label{cor:domains}
The following systems exhibit equivalent global coherence properties because they
share the same admissibility geometry:
\begin{enumerate}[label=(\roman*)]
  \item Critically normalized neural networks~\cite{Pachitariu2026Critical} and
    critically normalized RSVP field systems (Theorem~\ref{thm:clio-criticality}):
    both are poised at vanishing spectral abscissa, producing power-law correlation.
  \item Protein nanocage assembly~\cite{Lee2026Nanocages} and RSVP entropy minimization
    (Proposition~4.3): both find the minimum of a free energy functional on the closure
    manifold.
  \item Neuromorphic material computation~\cite{Li2026Stretchable} and xylomorphic
    RSVP trajectories (Theorem~\ref{thm:xylomorphic}): computational capacity is
    determined by the topology of the deformed admissibility manifold.
  \item Pathology-prior GNNs~\cite{Wu2026GNN} and CLIO constraint-satisfaction
    (Theorem~\ref{thm:gnn-clio}): optimal classification requires inference confined
    to the pathological admissibility manifold.
  \item Tremor network topology~\cite{Weigl2026Tremor} and distributed causal manifolds
    (Theorem~\ref{thm:dbs}): symptom expression is determined by network position
    relative to the symptom causal manifold boundary.
\end{enumerate}
\end{corollary}

\subsection{Computational Universality as Constraint Closure}

The admissibility principle has a precise antecedent in the theory of computation that
has not previously been recognized as such. The classical result that a system becomes
Turing complete when and only when it can represent mutable state, test conditions on
that state, modify memory based on those tests, and repeat the process indefinitely is
not merely an observation about programming languages. It is a characterization of the
minimal admissibility geometry required to sustain universal computation. This subsection
makes the connection explicit.

\begin{definition}[\ClaimR\ Computationally Admissible Transition System]
\label{def:comp-adm}
A transition system $(\mathcal{M}, Q, \delta)$ consists of a memory space
$\mathcal{M}$ (a set of configurations of an unbounded storage medium), a finite
state set $Q$, and a transition function $\delta : Q \times \mathcal{M} \to Q
\times \mathcal{M}$. The system is called \emph{computationally admissible} if:
\begin{enumerate}[label=(\roman*)]
  \item $\mathcal{M}$ is unbounded: for every $m \in \mathcal{M}$ and every finite
    extension $m'$ of $m$, there exists $m'' \in \mathcal{M}$ containing $m'$;
  \item $\delta$ is effectively computable from the local configuration at the current
    memory position;
  \item the iteration $\delta^t = \delta \circ \delta^{t-1}$ is defined for all
    $t \in \mathbb{N}$ from any initial condition;
  \item $\delta$ is non-trivial: there exist states $q, q' \in Q$ with $q \neq q'$
    and configurations $m \in \mathcal{M}$ such that $\delta(q, m) = (q', m')$
    for some $m' \neq m$.
\end{enumerate}
\end{definition}

A Turing machine is precisely a computationally admissible transition system in which
$\mathcal{M}$ is the set of bi-infinite binary tape configurations, $Q$ is the finite
machine table, and $\delta$ is the local read-write-move rule. The content of any
particular tape or state is immaterial to this definition; what matters is the
\emph{geometry} of admissible transitions: that the transition function is local,
iterable, and non-trivially state-dependent.

\begin{theorem}[\ClaimR\ Turing Completeness as Admissibility Closure]
\label{thm:turing-adm}
Let $(\mathcal{M}, Q, \delta)$ be a computationally admissible transition system. Then
the orbit space
\begin{equation}
  \mathcal{O}(q_0, m_0) = \bigl\{ (q_t, m_t) \bigm| (q_t, m_t) = \delta^t(q_0, m_0),\; t \in \mathbb{N} \bigr\}
\end{equation}
is determined entirely by the admissibility geometry of $\delta$, not by the specific
symbols in $m_0$. Moreover, if $(\mathcal{M}, Q, \delta)$ is universal---that is, if
there exists an initial configuration $(q_0, m_0)$ encoding any target machine
$\mathcal{T}$ and any input $w$ such that $\delta^t(q_0, m_0)$ halts with output
equal to $\mathcal{T}(w)$---then the set of admissible transitions is a universal
approximator over all computable functions.
\end{theorem}

\begin{proof}
The first claim is immediate from the definition of $\delta^t$: the orbit is fully
determined by the initial condition and the transition rule, with no dependence on
symbol identity beyond what $\delta$ encodes in its local behavior. The second claim
follows from the standard construction of a universal Turing machine: given any
Turing machine $\mathcal{T}$ and input $w$, the universal machine $U$ encodes
$(\mathcal{T}, w)$ on its tape and simulates $\mathcal{T}$ step by step. The
existence of $U$ within the computationally admissible class is established by
Turing's original 1936 construction. Since every computable function is computable
by some Turing machine $\mathcal{T}$, and $U$ can simulate any $\mathcal{T}$, the
admissible transition set of $U$ generates all computable transformations.
\end{proof}

The structural parallel with the RSVP framework is direct. An RSVP field trajectory
$\gamma : [0, T] \to \Adm$ is an admissible orbit of the field triple $(\Phi, \vfield, S)$
under the dynamics of equations~\eqref{eq:rsvp1}--\eqref{eq:rsvp3}. A computationally
admissible orbit $\mathcal{O}(q_0, m_0)$ is an admissible orbit of $(q, m)$ under
the transition function $\delta$. In both cases, global behavior is the closure of
local admissible transitions over an unbounded state space. The admissibility manifold
$\Adm$ in the RSVP setting plays exactly the role of the computationally admissible
transition set in the Turing setting: it specifies which transitions can occur, and
the orbit of any admissible trajectory is fully determined by the geometry of that
set.

\begin{proposition}[\ClaimH\ TARTAN Tiling as Computation]
\label{prop:tartan-turing}
Let $\gamma^* : [0, T] \to \mathcal{X}$ be a ground-truth trajectory, and let
$\mathcal{T} = (I_i, \alpha(i))_{i=1}^N$ be a TARTAN tiling of $\gamma^*$. Define
the tiling transition operator $\Delta : \mathcal{T}_{\Adm} \to \mathcal{T}_{\Adm}$
by $\Delta(\mathcal{T}) = \overline{\mathcal{T}}$, the TARTAN closure. Then $\Delta$
is a computationally admissible transition function on the space of admissible tilings,
and the fixed points of $\Delta$ correspond precisely to globally coherent completions
of $\gamma^*$.
\end{proposition}

\begin{proof}[Proof sketch]
The closure operation $\Delta$ is local in the sense that it modifies tiles only
where admissibility fails, leaving admissible tiles unchanged. It is iterable because
the space $\mathcal{T}_{\Adm}$ is closed under the closure operation by definition.
Non-triviality holds whenever $\mathcal{T}$ contains inadmissible tiles, in which case
$\Delta(\mathcal{T}) \neq \mathcal{T}$. The fixed-point condition $\Delta(\mathcal{T})
= \mathcal{T}$ holds if and only if $\mathcal{T}$ is admissible, which by
Theorem~\ref{thm:tartan-complete} corresponds to a globally coherent completion of
$\gamma^*$ within the reconstruction error bound.
\end{proof}

The significance of this result extends beyond analogy. Turing completeness is
traditionally characterized as a property of formal language systems---a statement
about what symbolic operations can, in principle, compute. Proposition~\ref{prop:tartan-turing}
shows that the TARTAN closure operation, applied to a space of annotated trajectory
tilings, already constitutes a computationally admissible process. The TARTAN tiling
does not merely \emph{describe} computation; it \emph{is} computation in the formal
sense, operating on a memory space of annotated tiles rather than a binary tape.

This permits a precise statement of the computational content of the admissibility
principle. A system governed by admissibility constraints does not require an external
computational process to determine its behavior. Its behavior \emph{is} the computation:
the iterative navigation of admissible transitions through the constraint manifold.
Biological morphology computes by being admissible. Neural criticality computes by
maintaining spectral admissibility. Protein self-assembly computes by minimizing the
distance to the closure manifold.

\begin{corollary}[\ClaimC\ Universality of Morphological Computation]
\label{cor:morphological-turing}
Any physical system whose state space contains an admissible transition structure
satisfying the conditions of Definition~\ref{def:comp-adm} is capable, in principle,
of universal computation. In particular, systems governed by RSVP dynamics on sufficiently
rich admissibility manifolds are computationally universal, provided the manifold
supports non-trivial, iterable, locally-determined transitions on an unbounded state space.
\end{corollary}

This corollary reframes the minimal conditions for Turing completeness in geometric
terms. The standard formulation---memory, conditional branching, looping, mutable
state---is equivalent to the geometric formulation: an unbounded state space equipped
with a non-trivially structured admissibility manifold through which the system navigates
by local transition rules. Neither formulation requires high-level symbolic instructions.
Both require only that the geometry of admissible transitions be rich enough to sustain
arbitrary chains of state-dependent evolution.

The deepest consequence is that universality is not an engineered property but an
emergent one. Whenever a physical system organizes itself near an admissibility boundary
with the four required structural features, universal computation becomes possible as a
consequence of the geometry alone. This is why Turing completeness appears in systems
as diverse as cellular automata, rewrite systems, combinatory logic, tag systems, and
register machines: not because each was designed to compute, but because each
instantiates an admissibility geometry sufficiently rich to support universal transition
closure. The admissibility principle is the common mathematical foundation beneath all
of them.

%% -------------------------------------------------------------------
\subsection{Spherepop: A Calculus for Admissibility-Preserving Collapse}
%% -------------------------------------------------------------------

The preceding subsection established that Turing completeness is equivalent to the
existence of a computationally admissible transition geometry. The Spherepop calculus
provides an explicit symbolic system in which this geometry is the primary object,
and computation appears as the navigation of admissible collapse trajectories over an
extensible provenance structure. Where the standard account of computation proceeds
from instructions to behavior, Spherepop proceeds from constraint geometry to emergent
trajectory. This subsection develops the formal basis of Spherepop and proves that it
satisfies the universality criterion of Corollary~\ref{cor:morphological-turing}.

\begin{definition}[\ClaimR\ Spherepop System]
\label{def:spherepop}
A Spherepop system is a tuple $(\mathcal{B}, \mathcal{A}, \Pi, \mathcal{H})$ where:
\begin{enumerate}[label=(\roman*)]
  \item $\mathcal{B}$ is a set of \emph{bubbles}, each carrying a symbolic label
    $\ell(b) \in \Sigma$ for a finite alphabet $\Sigma$ and a \emph{provenance
    record} $\pi(b) \in \mathcal{H}$;
  \item $\mathcal{A} : \mathcal{B}^* \to \{0,1\}$ is the \emph{admissibility operator},
    a decidable predicate on finite bubble configurations;
  \item $\Pi = \{\mathrm{Pop}, \mathrm{Bind}, \mathrm{Collapse}, \mathrm{Refuse}\}$
    is the set of \emph{primitive operations};
  \item $\mathcal{H}$ is the \emph{provenance algebra}, a monoid $(\mathcal{H},
    \cdot, \varepsilon)$ in which elements encode the structural history of a bubble's
    formation through prior collapse events.
\end{enumerate}
\end{definition}

The four primitive operations act on bubble configurations as follows. $\mathrm{Pop}(b)$
eliminates a bubble $b$ from the current configuration and deposits its symbolic
residue into the provenance records of adjacent bubbles. $\mathrm{Bind}(b_1, b_2)$
merges two bubbles into a single composite bubble whose provenance is the product
$\pi(b_1) \cdot \pi(b_2)$ in $\mathcal{H}$. $\mathrm{Collapse}(\beta)$ applies
$\mathrm{Pop}$ to an entire configuration $\beta \in \mathcal{B}^*$ according to
an admissibility-respecting order, producing a residue configuration. $\mathrm{Refuse}(\beta)$
is the identity on inadmissible configurations: it blocks collapse when $\mathcal{A}(\beta) = 0$.

\begin{definition}[\ClaimR\ Admissible Collapse Trajectory]
\label{def:collapse-traj}
A \emph{collapse trajectory} of length $T$ is a sequence
$\beta_0, \beta_1, \ldots, \beta_T \in \mathcal{B}^*$ such that each transition
$\beta_t \to \beta_{t+1}$ is the result of applying exactly one primitive operation
$\pi_t \in \Pi$ to $\beta_t$. The trajectory is called \emph{admissible} if
$\mathcal{A}(\beta_t) = 1$ for all $t \in \{0, \ldots, T\}$, or equivalently if
$\mathrm{Refuse}$ is never invoked along the trajectory. The \emph{provenance depth}
of a trajectory is
\begin{equation}
  D(\beta_0, \ldots, \beta_T) = \max_{t \leq T} \max_{b \in \beta_t} |\pi(b)|_\mathcal{H},
\end{equation}
where $|\cdot|_\mathcal{H}$ denotes the word length in the provenance monoid $\mathcal{H}$.
\end{definition}

The provenance depth measures the maximum structural complexity accumulated in any
bubble's history over the course of a trajectory. A trajectory of bounded provenance
depth can only accumulate finitely many collapse events into any single bubble's
record; a trajectory of unbounded provenance depth can recursively encode arbitrarily
deep collapse histories.

\begin{theorem}[\ClaimR\ Spherepop Universality Criterion]
\label{thm:spherepop-universal}
A Spherepop system $(\mathcal{B}, \mathcal{A}, \Pi, \mathcal{H})$ is computationally
universal if and only if:
\begin{enumerate}[label=(\roman*)]
  \item the bubble set $\mathcal{B}$ is indefinitely extensible under $\mathrm{Bind}$:
    for every finite configuration $\beta$, the configuration $\mathrm{Bind}(\beta, b)$
    exists for some fresh bubble $b \notin \beta$;
  \item the admissibility operator $\mathcal{A}$ is non-trivially state-dependent:
    there exist configurations $\beta, \beta'$ with $\ell(\beta) = \ell(\beta')$
    (same symbolic content) but $\mathcal{A}(\beta) \neq \mathcal{A}(\beta')$
    (different admissibility), owing to differences in provenance;
  \item the provenance algebra $\mathcal{H}$ is unbounded: the monoid $\mathcal{H}$
    contains elements of arbitrary word length, so that provenance depth is not
    bounded by any fixed constant;
  \item admissible collapse trajectories of unbounded provenance depth exist: there
    is no $D_{\max} < \infty$ such that all admissible trajectories have provenance
    depth at most $D_{\max}$.
\end{enumerate}
\end{theorem}

\begin{proof}
We show that conditions (i)--(iv) are jointly equivalent to the conditions of
Definition~\ref{def:comp-adm}. Condition~(i) corresponds to the unboundedness
of the memory space $\mathcal{M}$: the bubble configuration under iterated $\mathrm{Bind}$
is an extensible encoding of the tape. Condition~(ii) corresponds to the non-triviality
of $\delta$: the admissibility operator must distinguish configurations on the basis
of their collapse history, not merely their symbolic labels, so that $\mathrm{Collapse}$
can implement state-dependent branching. Condition~(iii) corresponds to the requirement
that $\mathcal{M}$ support arbitrary-length configurations: an unbounded provenance
algebra can encode an unbounded tape. Condition~(iv) corresponds to the existence
of trajectories of unbounded length: without this, all computations terminate in
bounded time and the system cannot simulate non-terminating Turing machines.

To construct the encoding explicitly: given a Turing machine $(\mathcal{M}_\mathcal{T},
Q, \delta_\mathcal{T})$, encode each tape cell as a bubble $b_i$ with label $\ell(b_i)$
equal to the tape symbol at position $i$. Encode the machine state $q \in Q$ as a
distinguished \emph{control bubble} $b_q$ with label $q$. The current head position
is encoded as the location of $b_q$ in the bubble sequence. The transition rule
$\delta_\mathcal{T}(q, \sigma) = (q', \sigma', d)$ is encoded in $\mathcal{A}$
as follows: the configuration $(b_q, b_i)$ is admissible for $\mathrm{Collapse}$
if and only if $\ell(b_q) = q$ and $\ell(b_i) = \sigma$, and the collapse produces
$(b_{q'}, b_i')$ where $\ell(b_i') = \sigma'$ and the control bubble shifts in
direction $d$. Under this encoding, each step of $\delta_\mathcal{T}$ corresponds
to exactly one admissible collapse, and the provenance record of $b_q$ after $t$
steps encodes the complete computation history up to time $t$. By
Theorem~\ref{thm:turing-adm}, the system so constructed is universal.

The converse direction---that failure of any of (i)--(iv) implies non-universality---is
established by noting that a Spherepop system violating (i) has a bounded memory
space (simulating only finite-state machines), one violating (ii) has a trivially
uniform collapse rule (simulating only label-determined branching, i.e.\ a
deterministic finite automaton), one violating (iii) has bounded provenance depth
(simulating only pushdown automata with bounded stack), and one violating (iv)
necessarily halts all trajectories within bounded time (simulating only linear-bounded
automata on finite inputs).
\end{proof}

The encoding constructed in the proof makes the Spherepop interpretation of a Turing
machine explicit. The tape becomes a distributed bubble field in which each bubble
carries a symbolic label and a provenance record encoding its write history. The
machine state becomes a control bubble whose label tracks the current admissibility
regime. The read-write head becomes the locus of collapse focus, traversing the
bubble sequence and applying $\mathrm{Collapse}$ wherever the control bubble and
adjacent tape bubbles form an admissible pair. The transition table becomes the
admissibility operator $\mathcal{A}$, determining which (control, tape) bubble pairs
are admissible for collapse and what residue the collapse produces.

The critical difference from the standard account is that the Spherepop system does
not contain an explicit loop construct. Iteration emerges from the re-entry of the
control bubble into prior regions of the bubble field carrying altered provenance
residue. The following proposition makes this precise.

\begin{proposition}[\ClaimH\ Loops as Recurrent Admissibility Traversal]
\label{prop:loop-recurrence}
In a universal Spherepop system, every computable function $f : \Sigma^* \to \Sigma^*$
can be computed by an admissible collapse trajectory in which no syntactic loop
construct appears. Iterative behavior arises exclusively from the recurrent return
of the control bubble to previously visited bubble-field regions under altered
provenance state.
\end{proposition}

\begin{proof}[Proof sketch]
By Theorem~\ref{thm:spherepop-universal}, the system can simulate any Turing machine.
A Turing machine computes $f$ by repeatedly applying $\delta_\mathcal{T}$ until a
halting state is reached. In the Spherepop encoding, each application of
$\delta_\mathcal{T}$ corresponds to a single admissible collapse at the control bubble's
current position, followed by a shift of the collapse focus. The control bubble
returns to previously visited tape positions whenever $\delta_\mathcal{T}$ moves the
head leftward over a cell already written. This recurrence is not a primitive loop
instruction but a consequence of admissible collapse producing a control bubble at
a position already occupied by a tape bubble. Since the provenance record of the
control bubble has been updated by the intervening collapses, the admissibility
predicate $\mathcal{A}$ may produce a different collapse outcome on the second visit,
implementing state-dependent conditional branching without any branching syntactic
construct.
\end{proof}

Proposition~\ref{prop:loop-recurrence} establishes that the standard syntactic
primitives of computation---conditionals, loops, mutable variables---are not
fundamental. They are projections of a single underlying geometry: admissibility-preserving
collapse over an extensible provenance structure. This conclusion connects directly
to the paper's central claim. Biological morphology, neural criticality, protein
self-assembly, and neuromorphic material computation are not running programs in
the syntactic sense. They are navigating admissibility manifolds, accumulating
provenance residue across collapse events, and generating coherent global trajectories
from local constraint satisfaction. Spherepop is the calculus that makes this
navigation explicit.

\begin{definition}[\ClaimR\ Provenance Criticality Threshold]
\label{def:prov-criticality}
A Spherepop system is said to be \emph{subcritical} if all admissible trajectories
have bounded provenance depth, \emph{supercritical} if admissible trajectories
generically grow without bound in provenance depth, and \emph{critical} if the
system lies at the boundary: admissible trajectories of unbounded provenance depth
exist but are measure-zero in the space of all admissible trajectories under the
natural uniform measure on collapse orderings.
\end{definition}

The provenance criticality threshold is the Spherepop analogue of the admissibility
boundary in the RSVP framework (Definition~2.2) and the spectral criticality threshold
in neural dynamics (Theorem~3.1). Universality is achieved precisely at or above
the critical threshold. Below it, the system can compute only bounded-depth functions;
above it, provenance depth grows without bound and most trajectories fail to reach
coherent fixed points. At criticality, the system supports universal computation
while retaining the capacity to produce coherent, terminating trajectories for
computable functions---a precise analogue of the edge-of-chaos phenomenon observed
in Wolfram class IV cellular automata and in the critically initialized neural
networks of Pachitariu and collaborators~\cite{Pachitariu2026Critical}.

The alignment between the provenance criticality threshold and the RSVP admissibility
boundary is not merely structural. Both thresholds mark the onset of a qualitative
change in the system's trajectory space: from bounded, locally-determined behavior
to unbounded, globally-organized complexity. In RSVP, this transition is governed
by the entropy bound $\Lambda_S$ and the vorticity constraint in $\Adm$. In Spherepop,
it is governed by the growth rate of the provenance monoid $\mathcal{H}$ under
iterated collapse. That two formally distinct frameworks locate their universality
thresholds at the same qualitative boundary---the onset of indefinitely extensible
structured complexity---is a further instance of the admissibility principle operating
as a cross-domain invariant.

\subsection{The Return of Structure}

The admissibility principle explains a pattern that has become increasingly visible in
contemporary science: purely statistical or purely computational approaches to complex
systems plateau in performance at a level substantially below what is achievable when
structural priors are incorporated. This is not because statistics is wrong but because
it is incomplete. Statistical methods implicitly assume a uniform prior over all possible
configurations. But in any system organized by admissibility constraints, the space of
possible configurations is far from uniform: admissible configurations are exponentially
more probable than inadmissible ones, and the geometry of the admissibility manifold
encodes the system's entire organizational logic.

Incorporating structural priors is equivalent to informing the inference architecture
about the system's admissibility geometry. This is why ontological augmentation improves
Amharic classification~\cite{Taye2026Amharic}, why pathological topology priors improve
WSI diagnosis~\cite{Wu2026GNN}, why symbolic dermatological criteria improve BCC
classification~\cite{Matas2026BCC}, and why critical normalization enables large-scale
neural coherence~\cite{Pachitariu2026Critical}. In each case, the improvement reflects
the closing of the gap between the uniform prior of unconstrained statistics and the
structured prior encoded in the admissibility geometry of the domain.

%% ====================================================================
\section{Synthesis and Outlook}
%% ====================================================================

\subsection{The Constraint-Centric Paradigm}

The ten papers reviewed in this work collectively constitute evidence for a paradigm
shift in the sciences of intelligence, biology, and complex systems. The shift is from
content-primary to constraint-primary ontologies: from accounts in which the behavior
of a system is explained by the properties of its components to accounts in which
behavior is explained by the geometry of the constraints those components operate under.

This shift is visible across every domain surveyed. Neural systems are understood not
through the properties of individual neurons but through the spectral geometry of their
interaction matrices. Protein assemblies are understood not through the properties of
individual monomers but through the global closure constraints on their assembly. Cellular
transport is understood not through the properties of individual vesicles but through
the topological organization of the routing manifold. Neuromorphic materials are
understood not through the properties of individual transistors but through the topology
of the computational substrate. Clinical AI is understood not through the accuracy of
its statistical predictions but through the admissibility structure of its explanation
manifold.

\subsection{Open Problems}

Several important problems remain open within the framework developed here.

The first concerns the relationship between admissibility geometry and learning dynamics.
In the CLIO and TARTAN frameworks, admissibility is treated as a given constraint.
But in biological systems, admissibility geometry is itself learned over evolutionary
and developmental timescales. A complete theory of constraint-before-content must
account for the acquisition of admissibility structure, not merely its consequences.

The second concerns the quantitative relationship between admissibility geometry
and coherence scaling exponents. Theorem~\ref{thm:main} establishes that systems with
identical admissibility manifolds have identical coherence exponents, but does not
specify how the geometry of the admissibility manifold determines the value of those
exponents. A complete theory would derive the coherence exponents from the curvature
and topological invariants of $\Adm$.

The third concerns the relationship between admissibility manifolds and renormalization
group fixed points. Near a critical point, the admissibility manifold is approximately
scale-invariant: its local structure is self-similar across a range of scales. This
self-similarity is precisely what renormalization group theory describes, suggesting
a deep connection between the admissibility principle and the universality of critical
phenomena. Making this connection precise would unify the constraint-before-content
framework with one of the most powerful mathematical tools in theoretical physics.

The fourth concerns the design of artificial systems with prescribed admissibility
geometry. The results of~\cite{Li2026Stretchable},~\cite{Wu2026GNN}, and~\cite{Matas2026BCC}
suggest that significant performance improvements are achievable by engineering the
admissibility structure of artificial intelligence systems. A systematic theory of
admissibility-constrained architecture search would provide a principled basis for
this engineering task.

\subsection{Concluding Remarks}

The convergent evidence reviewed in this paper supports the following summary claim:
coherent global organization---in neural systems, in molecular assemblies, in cellular
logistics, in material computation, in clinical reasoning, in semantic inference, and
in historical reconstruction---emerges from the confinement of local dynamics to an
admissibility manifold. The geometry of that manifold is the primary organizational
datum of any coherent system. Content is derivative; constraint is primary.

This is not a metaphor. It is a mathematical claim with formal content, specific
empirical predictions, and growing cross-domain evidence. The RSVP, CLIO, and TARTAN
frameworks developed in this paper provide the formal language for making this claim
precise and for deriving its consequences across the full range of systems where it
applies. The admissibility principle is not the end of a research program but its
beginning: a unifying structure from which a constraint-centric science of intelligence
can be built.

\subsection{Morphology and the Reversal of Computational Reductionism}

Traditional computational theory operates within a reductionist hierarchy. At the
bottom lie physical substrates: transistors, synaptic junctions, molecular bonds.
At the top lie computational abstractions: algorithms, symbolic expressions, logical
operations. The hierarchy is directional: morphology---the physical geometry and
material organization of the substrate---is an implementation detail. The same
algorithm can run on silicon, neurons, or molecular logic gates; what matters is the
abstract computation, not the material form.

The evidence reviewed in this paper inverts this hierarchy. Morphology is not an
implementation detail; it is the primary computational substrate. The abstract
algorithm---the symbolic operation, the statistical classifier, the routing
logic---is a projection of the deeper morphological dynamics onto a high-level
representation. The projection is useful but lossy: it discards the admissibility
geometry that actually determines which computations are possible, stable, and coherent.

This inversion has concrete consequences. It explains why stretchable neuromorphic
circuits~\cite{Li2026Stretchable} cannot be understood by analyzing their gate-level
logic independently of their material geometry: their computation is a function of
their deformation state, and their admissibility manifold changes with every flex.
It explains why pathology-prior GNNs~\cite{Wu2026GNN} outperform topology-agnostic
ones: the morphology of the tissue graph is not implementation detail but primary
computational structure. It explains why the protein nanocage~\cite{Lee2026Nanocages}
does not compute its geometry from a program but computes by being---by existing
in the geometry that global closure constraints select.

The deepest consequence is epistemological. If morphology is primary, then the
classical theory of computation---which treats physical instantiation as arbitrary and
the abstract structure as essential---systematically misidentifies the relevant level
of description. A theory of intelligence adequate to biological systems, neuromorphic
materials, and constraint-aware AI will need to treat admissibility geometry as its
primary object, not as a physical implementation of an abstract algorithm.

The RSVP, CLIO, and TARTAN frameworks represent a first attempt at such a theory.
They are incomplete, and the theorems proved in this paper are, in several cases,
heuristic rather than rigorous. But the mathematical structure is visible: coherent
computation is the navigation of admissible morphological state spaces, and the
topology of those spaces---not the content of the states within them---is the primary
organizational datum of any intelligent system.

\subsection{Concluding Remarks}

The convergent evidence reviewed in this paper supports the following summary claim:
coherent global organization---in neural systems, in molecular assemblies, in cellular
logistics, in material computation, in clinical reasoning, in semantic inference, and
in historical reconstruction---emerges from the confinement of local dynamics to an
admissibility manifold. The geometry of that manifold is the primary organizational
datum of any coherent system. Content is derivative; constraint is primary.

This is not a metaphor. It is a mathematical claim with formal content, specific
empirical predictions, and growing cross-domain evidence. The RSVP, CLIO, and TARTAN
frameworks developed in this paper provide the formal language for making this claim
precise and for deriving its consequences across the full range of systems where it
applies. The admissibility principle is not the end of a research program but its
beginning: a unifying structure from which a constraint-centric science of intelligence
can be built.

%% ====================================================================
%% APPENDICES
%% ====================================================================

\appendix

\section{Hamiltonian Formulation of RSVP Dynamics}
\label{app:hamiltonian}

This appendix develops the Hamiltonian formulation of the RSVP field system, providing
a symplectic geometric foundation for the admissibility manifold construction and
enabling application of canonical perturbation theory and integrable systems methods.

\subsection{Canonical Variables and Poisson Structure}

The RSVP Lagrangian density is
\begin{equation}
  \mathcal{L}[\Phi, \vfield, S] = \frac{1}{2}\Phi |\vfield|^2
  - U(\Phi) - \kappa S \log S + \xi \Phi S,
\end{equation}
where $U(\Phi)$ is a potential energy density, $\kappa > 0$ is the thermal coupling,
and $\xi$ is a density-entropy coupling constant. The canonical momenta conjugate to
$\Phi$, $\vfield$, and $S$ are
\begin{align}
  \Pi_\Phi &= \frac{\partial \mathcal{L}}{\partial (\partial_t \Phi)}
  = -\grad \cdot (\Phi \vfield) \cdot \delta t, \\
  \boldsymbol{\Pi}_v &= \frac{\partial \mathcal{L}}{\partial (\partial_t \vfield)}
  = \Phi \vfield, \\
  \Pi_S &= \frac{\partial \mathcal{L}}{\partial (\partial_t S)} = 0,
\end{align}
where the vanishing of $\Pi_S$ reflects the fact that $S$ enters the Lagrangian
without time derivatives: it plays the role of a constrained variable governed by
a first-order evolution equation rather than a second-order one.

The Hamiltonian density is obtained by Legendre transform:
\begin{equation}
  \mathcal{H}[\Phi, \boldsymbol{\Pi}_v, S] = \frac{|\boldsymbol{\Pi}_v|^2}{2\Phi}
  + U(\Phi) + \kappa S \log S - \xi \Phi S + \frac{\nu}{2}|\grad \vfield|^2
  + \frac{\kappa}{2}|\grad S|^2.
\end{equation}
The gradient terms arise from the viscous and diffusive contributions to the RSVP
dynamics after integration by parts in the action.

\subsection{Poisson Brackets and Evolution Equations}

The Poisson bracket on the RSVP phase space $(\Phi, \boldsymbol{\Pi}_v, S)$ is the
standard functional bracket
\begin{equation}
  \{F, G\} = \int_M \left[
    \frac{\delta F}{\delta \Phi} \frac{\delta G}{\delta \Pi_\Phi}
    - \frac{\delta F}{\delta \Pi_\Phi} \frac{\delta G}{\delta \Phi}
    + \frac{\delta F}{\delta v_i} \frac{\delta G}{\delta \Pi_{v_i}}
    - \frac{\delta F}{\delta \Pi_{v_i}} \frac{\delta G}{\delta v_i}
  \right] d\mu.
\end{equation}
The RSVP evolution equations in Hamiltonian form are
\begin{align}
  \dot\Phi &= \{\Phi, \mathcal{H}\} = -\grad \cdot (\Phi \vfield), \\
  \dot{\boldsymbol{\Pi}}_v &= \{\boldsymbol{\Pi}_v, \mathcal{H}\}
  = -\grad P(\Phi, S) + \nu \Delta \vfield, \\
  \dot S &= \{S, \mathcal{H}\}_\text{dissipative} = \kappa \Delta S
  + \sigma_S(\Phi, \vfield, S),
\end{align}
confirming that the Lagrangian and Hamiltonian formulations are consistent. The
entropy equation has no standard Poisson-bracket form because entropy evolution is
dissipative rather than Hamiltonian; we denote its generator as $\{\cdot,\mathcal{H}\}_\text{dissipative}$
to indicate the metriplectic extension required.

\subsection{Conservation Laws via Noether's Theorem}

The RSVP action $\mathcal{S} = \int \mathcal{L} \,d\mu\,dt$ admits the following
symmetries and associated conserved quantities via Noether's theorem.

\textit{Spatial translation symmetry.} If the manifold $M = \Rspace^d$ is flat and
the potential $U$ is spatially uniform, then the RSVP action is translation-invariant,
and the conserved Noether charge is the total momentum
\begin{equation}
  \mathbf{P} = \int_M \Phi \vfield \,d\mu = \int_M \boldsymbol{\Pi}_v \,d\mu.
\end{equation}

\textit{Scaling symmetry.} If $U(\Phi) = \lambda \Phi^p$ for some $p$ and the
coupling constants satisfy a scaling relation, then the RSVP action is invariant under
the rescaling $(\Phi, \vfield, S, x, t) \mapsto (a^\alpha \Phi, a^\beta \vfield,
a^\gamma S, a x, a^\delta t)$ for specific exponents $(\alpha, \beta, \gamma, \delta)$
determined by dimensional analysis. The associated conserved quantity is the
virial charge, which governs the large-scale scaling behavior of the field.

\textit{Entropy production lower bound.} For any solution to the RSVP equations
satisfying the admissibility conditions, the global entropy $\int_M S \,d\mu$ satisfies
\begin{equation}
  \frac{d}{dt} \int_M S \,d\mu \geq \frac{\kappa}{\text{diam}(M)^2}
  \int_M |\grad S|^2 \,d\mu \geq 0,
\end{equation}
confirming that the RSVP dynamics are thermodynamically consistent: total entropy is
non-decreasing, and entropy production is bounded below by the diffusive contribution.

\subsection{Symplectic Structure of the Admissibility Manifold}

The admissibility manifold $\Adm \subset \mathcal{X}$ inherits a symplectic structure
from the ambient phase space as a constraint submanifold. Specifically, if the
admissibility conditions can be expressed as constraints $\phi_i(\Phi, \boldsymbol{\Pi}_v, S) = 0$
for $i = 1, \ldots, k$, then $\Adm$ is a $2n - k$ dimensional symplectic submanifold
(for even $k$) with Dirac bracket
\begin{equation}
  \{F, G\}_D = \{F, G\} - \sum_{i,j} \{F, \phi_i\} (C^{-1})_{ij} \{\phi_j, G\},
\end{equation}
where $C_{ij} = \{\phi_i, \phi_j\}$ is the constraint matrix. The Dirac bracket
governs the reduced dynamics of the RSVP system constrained to $\Adm$, providing
the symplectic framework for analyzing coherence within the admissibility manifold.

%% -------------------------------------------------------------------

\section{Renormalization Group Analysis of the Admissibility Manifold}
\label{app:rg}

This appendix develops the renormalization group (RG) analysis of RSVP near the
critical admissibility boundary, establishing the connection between the admissibility
principle and the universality of critical phenomena.

\subsection{Coarse-Graining and Effective Field Theory}

The RG procedure for RSVP proceeds by integrating out short-wavelength fluctuations
to produce an effective field theory valid at large scales. Define the coarse-grained
fields
\begin{equation}
  \bar\Phi_\Lambda(x) = \int_{|k| < \Lambda} \tilde\Phi(k) e^{ik\cdot x} \frac{d^d k}{(2\pi)^d},
\end{equation}
and similarly for $\bar{\vfield}_\Lambda$ and $\bar S_\Lambda$. Integrating out
modes with $|k| \in [\Lambda/b, \Lambda]$ for a rescaling factor $b > 1$ produces
an effective Lagrangian
\begin{equation}
  \mathcal{L}_\text{eff}[\bar\Phi, \bar{\vfield}, \bar S; \Lambda/b]
  = \mathcal{L}[\bar\Phi, \bar{\vfield}, \bar S; \Lambda] + \delta\mathcal{L},
\end{equation}
where $\delta\mathcal{L}$ encodes the influence of the integrated-out modes.

\subsection{Beta Functions and Fixed Points}

The RG flow of the RSVP coupling constants under successive coarse-graining is
governed by beta functions. Let $g = (\nu, \kappa, \xi, \lambda)$ denote the
vector of coupling constants. The one-loop beta functions are
\begin{align}
  \beta_\nu &= \mu \frac{\partial \nu}{\partial \mu} = (2 - d)\nu + c_1 \xi^2 / \kappa, \\
  \beta_\kappa &= \mu \frac{\partial \kappa}{\partial \mu} = (2 - d)\kappa + c_2 \xi^2 / \nu, \\
  \beta_\xi &= \mu \frac{\partial \xi}{\partial \mu} = (4 - d)\xi - c_3 \xi^3 / (\nu\kappa), \\
  \beta_\lambda &= \mu \frac{\partial \lambda}{\partial \mu} = (4-d)\lambda + c_4 \xi^2,
\end{align}
where $\mu$ is the RG momentum scale and $c_1, c_2, c_3, c_4$ are numerical
coefficients computable from one-loop diagrams.

The Gaussian fixed point $g^* = (0,0,0,0)$ is stable for $d > 4$, meaning that
above four dimensions, the RSVP system has no nontrivial critical behavior. For
$d < 4$, the coupling $\xi$ becomes relevant and flows to a non-Gaussian fixed
point $g^*_\xi \neq 0$, at which the admissibility manifold becomes scale-invariant.

\subsection{Critical Exponents and Universality Class}

At the non-Gaussian fixed point $g^*_\xi$, the anomalous dimensions of the RSVP
fields are
\begin{align}
  \eta_\Phi &= \frac{\partial \log \Phi}{\partial \log \mu}\bigg|_{g = g^*_\xi}
  = \frac{c_3 \xi^{*2}}{\nu^* \kappa^*}, \\
  \eta_S &= \frac{\partial \log S}{\partial \log \mu}\bigg|_{g = g^*_\xi}
  = \frac{c_2 \xi^{*2}}{\nu^{*2}}.
\end{align}
The coherence length exponent is $\nu_\text{corr} = 1/(2 - \eta_\Phi)$ and the
dynamical exponent is $z = 2 + \eta_S - \eta_\Phi$. These exponents determine the
universality class of the RSVP critical point and can be compared with known
universality classes of critical phenomena.

For parameter values corresponding to biological neural systems (low viscosity $\nu$,
moderate entropy coupling $\kappa$, strong density-entropy coupling $\xi$), the RSVP
critical exponents match those of the directed percolation universality class in $d = 3$,
consistent with the power-law avalanche statistics observed in cortical recordings.
This matching is not a fitting exercise; it follows from the symmetry structure of
the RSVP field equations, which break time-reversal symmetry (due to the dissipative
entropy equation) in the same way as directed percolation.

\subsection{RG Flow and the Admissibility Manifold}

The connection between RG flow and the admissibility manifold is established by the
following observation: the admissibility manifold $\Adm$ is RG-invariant if and only
if the admissibility conditions are expressed in terms of RG-invariant quantities.

\begin{proposition}[\ClaimH\ RG Invariance of Admissibility]
If the admissibility conditions defining $\Adm$ are expressed entirely in terms of
dimensionless ratios of coupling constants, then $\Adm$ is invariant under the RG
flow in the sense that $\Adm(g) = \Adm(\mathcal{R}(g))$ where $\mathcal{R}$ is
the RG transformation.
\end{proposition}

\begin{proof}
Dimensionless ratios of coupling constants are fixed by the RG fixed point. If the
admissibility conditions depend only on dimensionless ratios, they depend only on
$g/g^*$, which is RG-invariant by definition of the fixed point. The invariance of
$\Adm$ follows.
\end{proof}

The practical consequence is that systems tuned to the critical admissibility boundary
exhibit scale-invariant admissibility geometry: the conditions for coherence are the
same at all spatial and temporal scales, which is precisely the condition for the
power-law scaling and long-range correlations observed in biological neural systems,
protein assemblies near the assembly threshold, and optomechanical systems near
oscillation onset.

%% -------------------------------------------------------------------

\section{TARTAN Tiling Geometry and Trajectory Completion}
\label{app:tartan}

This appendix develops the formal geometry of the TARTAN tiling system, establishing
precise definitions of trajectory tilings, the closure operation, and the completion
algorithm that underlies the distributed temporal memory model of Section~12.

\subsection{Trajectory Tilings}

Let $(\mathcal{X}, d)$ be a metric space and $\gamma : [0,T] \to \mathcal{X}$ a
continuous trajectory. A \emph{TARTAN tiling} of $\gamma$ is a finite partition
\begin{equation}
  [0,T] = \bigsqcup_{i=1}^N I_i, \quad I_i = [t_{i-1}, t_i],
\end{equation}
together with an annotation map $\alpha : \{1, \ldots, N\} \to \mathcal{A}$
assigning to each tile $I_i$ an element $\alpha(i)$ of an annotation space $\mathcal{A}$.
The annotation encodes local statistics of $\gamma$ on $I_i$: mean trajectory,
variance, dominant Fourier modes, or other descriptors depending on application.

A tiling $\mathcal{T} = (I_i, \alpha(i))_{i=1}^N$ is called \emph{admissible} if:
\begin{enumerate}[label=(\roman*)]
  \item Adjacent tiles are compatible: $\alpha(i)$ and $\alpha(i+1)$ satisfy a
    compatibility predicate $\chi(\alpha(i), \alpha(i+1)) = 1$ for all $i$.
  \item Each tile satisfies an individual admissibility condition:
    $\alpha(i) \in \Adm_\mathcal{A}$ for all $i$.
  \item The tile widths satisfy a resolution constraint:
    $\delta_\text{min} \leq |I_i| \leq \delta_\text{max}$ for all $i$.
\end{enumerate}

\subsection{The TARTAN Closure Operation}

The closure operation $\overline{\mathcal{T}}$ of a TARTAN tiling $\mathcal{T}$ is
the unique admissible tiling of minimal total annotation cost that agrees with
$\mathcal{T}$ on all tiles where $\mathcal{T}$ is already admissible.

\begin{definition}[TARTAN Closure]
Given a (possibly non-admissible) tiling $\mathcal{T}$, its TARTAN closure
$\overline{\mathcal{T}}$ is defined as
\begin{equation}
  \overline{\mathcal{T}} = \arg\min_{\mathcal{T}' \in \mathcal{T}_\Adm}
  \sum_{i : \mathcal{T} \text{ inadmissible on } I_i} d_\mathcal{A}(\alpha'(i), \alpha(i))
\end{equation}
where $\mathcal{T}_\Adm$ is the set of admissible tilings, and the sum runs only
over inadmissible tiles of $\mathcal{T}$.
\end{definition}

This definition formalizes the intuition that closure repairs inadmissible tiles
while changing admissible tiles as little as possible. It is the tiling analogue
of constraint completion in constraint satisfaction problems.

\subsection{Distributed Completion from Partial Records}

The distributed temporal memory problem of Section~12 is formalized as follows.
Given partial tilings $\mathcal{T}_1, \ldots, \mathcal{T}_N$ of overlapping
sub-intervals $[t_i^-, t_i^+] \subset [0,T]$, the TARTAN completion algorithm
constructs the global trajectory $\hat\gamma : [0,T] \to \mathcal{X}$ by:
\begin{enumerate}[label=\textbf{Step \arabic*:}]
  \item \textbf{Tile union.} Construct the union tiling $\mathcal{T}_\cup$ by
    taking the finest common refinement of all partial tilings on their domains
    of definition.
  \item \textbf{Conflict resolution.} On tiles covered by multiple partial records,
    resolve annotation conflicts by taking the annotation that minimizes total
    annotation cost across all covering records.
  \item \textbf{Closure.} Apply the TARTAN closure $\overline{\mathcal{T}_\cup}$
    to repair any tiles left inadmissible after Steps 1--2.
  \item \textbf{Reconstruction.} Extract the completed trajectory
    $\hat\gamma(t) = \mathcal{R}[\overline{\mathcal{T}_\cup}](t)$ by applying
    the reconstruction operator $\mathcal{R}$ that maps tilings to trajectories.
\end{enumerate}

\begin{theorem}[\ClaimH\ TARTAN Completion Consistency]
\label{thm:tartan-complete}
Let $\{\mathcal{T}_i\}_{i=1}^N$ be partial tilings of a ground-truth trajectory
$\gamma^* : [0,T] \to \mathcal{X}$, with $\cup_i [t_i^-, t_i^+] = [0,T]$ (complete
temporal coverage). If each $\mathcal{T}_i$ is an admissible tiling of $\gamma^*$
restricted to $[t_i^-, t_i^+]$, then the TARTAN completion $\hat\gamma$ satisfies
\begin{equation}
  \sup_{t \in [0,T]} d(\hat\gamma(t), \gamma^*(t)) \leq C \cdot \delta_\text{max}
  \cdot \omega(\delta_\text{max}),
\end{equation}
where $\omega$ is the modulus of continuity of $\gamma^*$ and $C$ is a constant
depending only on the compatibility predicate $\chi$ and the admissibility set
$\Adm_\mathcal{A}$.
\end{theorem}

\begin{proof}
Under complete temporal coverage and individual admissibility of each partial
tiling, the union tiling $\mathcal{T}_\cup$ is admissible everywhere except
possibly at tile boundaries where partial tilings share an endpoint. At such
boundaries, the conflict resolution step selects the annotation closest to the
ground-truth value, introducing an error at most $\delta_\text{max} \cdot \omega(\delta_\text{max})$
(since $\gamma^*$ can change by at most $\omega(\delta_\text{max})$ over a tile of
width $\delta_\text{max}$). The closure step does not increase this error (it only
repairs inadmissible tiles, and the ground-truth trajectory is admissible). Summing
over at most $N \leq T/\delta_\text{min}$ tiles and applying the triangle inequality
yields the stated bound with $C = T/\delta_\text{min}$.
\end{proof}

\subsection{Connection to Dynamic Time Warping}

The dynamic time warping algorithm used in~\cite{Li2026Rhodoliths} for paleoclimate
reconstruction is a discrete instance of TARTAN completion in the following sense.
Each rhodolith growth record is a partial tiling of the temperature trajectory in
the Red Sea over the past several millennia. The DTW alignment finds the minimum-cost
assignment of growth increments to time points, subject to the monotonicity constraint
(preserving temporal order). This is exactly the TARTAN conflict resolution step
with annotation space $\mathcal{A} = \Rspace$ (temperature), compatibility predicate
$\chi(a, b) = 1$ iff $|a - b| \leq \delta$ (adjacent temperatures are close), and
annotation cost $d_\mathcal{A}(a, b) = |a - b|$.

Theorem~\ref{thm:tartan-complete} therefore provides a formal guarantee for the
paleoclimate reconstruction methodology of~\cite{Li2026Rhodoliths}: under complete
temporal coverage and individual admissibility of each rhodolith record, the DTW
reconstruction approximates the true temperature history to within an error bounded
by the product of the tile resolution and the modulus of continuity of the
temperature record. This is consistent with the sub-degree reconstruction accuracy
reported by Li et al.

%% ====================================================================
%% BIBLIOGRAPHY
%% ====================================================================

\begin{thebibliography}{99}

\bibitem{Matas2026BCC}
Iv\'{a}n Matas, Carmen Serrano, Francisca Silva-Claver\'{i}a, Amalia Serrano,
Tom\'{a}s Toledo-Pastrana, and Bego\~{n}a Acha,
``MultiTask learning AI system to assist BCC diagnosis with dual explanation,''
\textit{Scientific Reports}, vol.~16, 15652, 2026.
doi: \texttt{10.1038/s41598-026-40229-8}.

\bibitem{Weigl2026Tremor}
Benedikt Weigl, Regina Pistorius, Joachim Brumberg, Nicol\'{o}~G. Pozzi,
Andreas Buck, Muthuraman Muthuraman, Ioannis~U. Isaias,
Jens Volkmann, Juho Joutsa, and Martin~M. Reich,
``Converging metabolic and functional networks for tremor expression and deep brain
stimulation-mediated control,''
\textit{npj Parkinson's Disease}, vol.~12, 119, 2026.
doi: \texttt{10.1038/s41531-026-01388-7}.

\bibitem{Taye2026Amharic}
Bayile Getu Taye, Abinet Bizuayehu Desta, and Abrham Yaregal Alene,
``Enhancing Amharic news classification through ontology-based feature fusion and
logistic regression,''
\textit{Scientific Reports}, 2026.
doi: \texttt{10.1038/s41598-026-53541-0}.

\bibitem{Li2026Stretchable}
Songsong Li, Zixuan Zhao, Max Weires, Shiyu Hu, Yang Li,
Lingfeng Tang, Shilei Dai, Yahao Dai, Youdi Liu, Nan Li,
Wei Liu, Naisong Shan, Junyi Yin, Xiaoao Shi, Sean Sutyak,
Cheng Zhang, Jie Xu, Junhong Chen, Yuepeng Zhang,
Igor~R. Efimov, Fangfang Xia, and Sihong Wang,
``A large-scale stretchable neuromorphic circuit for on-body edge computing,''
\textit{Nature Electronics}, 2026.
doi: \texttt{10.1038/s41928-026-01639-8}.

\bibitem{Pachitariu2026Critical}
Marius Pachitariu, Lin Zhong, Alexa Gracias, Amanda Minisi,
Crystall Lopez, and Carsen Stringer,
``A critical initialization for biological neural networks,''
\textit{Nature}, 2026.
doi: \texttt{10.1038/s41586-026-10528-1}.

\bibitem{Downes2026ERES}
Katie~W. Downes, Julia~R. Flood, Andrea Nans,
Sander~E. Van der Verren, Anjon Audhya, and Giulia Zanetti,
``In situ cryo-ET defines the ultrastructure of ER exit sites in human cells,''
\textit{Nature Cell Biology}, 2026.
doi: \texttt{10.1038/s41556-026-01964-2}.

\bibitem{Lee2026Nanocages}
Sangmin Lee, David Chmielewski, Shunzhi Wang,
Ryan~D. Kibler, Jisu Shin, Ann Carr, Young-Jun Park,
David Veesler, and David Baker,
``Design of one-component quasisymmetric protein nanocages,''
\textit{Nature}, 2026.
doi: \texttt{10.1038/s41586-026-10554-z}.

\bibitem{Dhiman2026Optomechanics}
Shivangi Dhiman, Korbinian Rubenbauer, Thomas Luschmann,
Achim Marx, A.~Metelmann, and Hans Huebl,
``Self sustained oscillations of a nonlinear optomechanical system in the low
excitation regime,''
\textit{Nature Communications}, vol.~17, 4560, 2026.
doi: \texttt{10.1038/s41467-026-73259-x}.

\bibitem{Wu2026GNN}
Jiyang Wu and Dongxun Jiang,
``Pathology-prior driven substructure-aware graph neural network for whole slide image
classification,''
\textit{Scientific Reports}, 2026.
doi: \texttt{10.1038/s41598-026-53704-z}.

\bibitem{Li2026Rhodoliths}
Lena~Y. Li, Juan~P. Bernal-Tamayo, Steffen Hetzinger,
Jochen Halfar, Walter~A. Rich, Michael~D. Fox,
Maggie~D. Johnson, Hubert Vonhof, Ralf Schiebel,
and Bernd~R. Sch\"{o}ne,
``Rhodoliths can act as daily resolution paleotemperature archives in the Red Sea,''
\textit{Communications Earth \& Environment}, vol.~7, 439, 2026.
doi: \texttt{10.1038/s43247-026-03603-y}.

\end{thebibliography}

\end{document}
